\documentclass[11pt]{article}
\usepackage[a4paper,margin=28mm]{geometry}
\usepackage[T1]{fontenc}
\usepackage{lmodern}
\usepackage{amsmath,amssymb,amsthm,mathtools}
\usepackage{microtype}
\usepackage{enumitem}
\usepackage{needspace}
\usepackage[hidelinks]{hyperref}
\hypersetup{pdftitle={Nuclear dimension and Jiang-Su stability without elementary subquotients},
  pdfauthor={OpenAI},
  pdfsubject={Nuclear dimension, Jiang-Su stability, and full signs in maximal tensor powers}}
\numberwithin{equation}{section}
\newtheorem{theorem}{Theorem}[section]
\newtheorem{lemma}[theorem]{Lemma}
\newtheorem{proposition}[theorem]{Proposition}
\newtheorem{corollary}[theorem]{Corollary}
\theoremstyle{definition}
\newtheorem{definition}[theorem]{Definition}
\theoremstyle{remark}
\newtheorem{remark}[theorem]{Remark}
\newcommand{\Z}{\mathcal Z}
\newcommand{\Cu}{\mathrm{Cu}}
\newcommand{\nucdim}{\dim_{\mathrm{nuc}}}
\newcommand{\Ideal}{\operatorname{Ideal}}
\newcommand{\her}{\operatorname{her}}
\newcommand{\N}{\mathbb N}
\newcommand{\C}{\mathbb C}
\newcommand{\R}{\mathbb R}
\newcommand{\Q}{\mathbb Q}
\newcommand{\K}{\mathcal K}
\newcommand{\spc}{\operatorname{sp}}
\newcommand{\ann}{\operatorname{Ann}}
\newcommand{\id}{\mathrm{id}}
\DeclareMathOperator{\diag}{diag}
\DeclareMathOperator{\dist}{dist}
\DeclareMathOperator{\supp}{supp}
\DeclareMathOperator{\Ad}{Ad}
\setlist[enumerate]{itemsep=3pt,topsep=5pt}
\setlist[itemize]{itemsep=3pt,topsep=5pt}
\raggedbottom
\title{Nuclear dimension and Jiang--Su stability\\without elementary subquotients}
\author{OpenAI}
\date{September 23, 2026}
\begin{document}
\maketitle
\begin{abstract}
For separable nuclear $C^*$-algebras with no nonzero elementary ideal
subquotients, finite nuclear dimension is equivalent to Jiang--Su
stability. More generally, $\dim_{\mathrm{nuc}}(A_0\otimes\mathcal Z)\le1$
for every separable nuclear $A_0$. This proves Robert and Tikuisis's
Conjecture (C1) and the nuclear-dimension equivalence in the nonsimple
Toms--Winter regularity question.
\end{abstract}
\section{Introduction}
\label{sec:introduction}

Nuclear dimension and Jiang--Su absorption describe two forms of
regularity in a nuclear $C^*$-algebra. Nuclear dimension counts the
order-zero pieces needed in finite-dimensional completely positive
approximations; an order-zero map preserves orthogonality of positive
elements. Jiang--Su stability means an isomorphism
$A\cong A\otimes\Z$, where $\Z$ is the algebra introduced by Jiang and Su \cite{JS}. The first
condition gives a bound on approximation complexity, while the second
expresses invariance under tensoring with a fixed infinite-dimensional
algebra. Their relationship connects noncommutative dimension with
the structure and classification of nuclear $C^*$-algebras.

The obstruction to passing from finite dimension to absorption is
already visible in a matrix algebra: it has nuclear dimension zero
but does not absorb $\Z$. The relevant exclusion must apply throughout
the ideal structure. A $C^*$-algebra $A$ is \emph{nowhere scattered}
if no quotient $I/J$, for closed two-sided ideals $J\subseteq I\subseteq A$,
is a nonzero elementary algebra $\mathcal K(H)$. This is equivalent
to the absence of finite-dimensional irreducible representations in
every hereditary subalgebra \cite[Theorem~3.1]{TV}.

\begin{theorem}[Nonsimple nuclear regularity]
\label{thm:main}
\label{intro:regularity}
Let $A$ be a separable nuclear complex $C^*$-algebra with no nonzero
elementary ideal subquotients. The following conditions are equivalent:
\begin{enumerate}
\item $\nucdim(A)<\infty$;
\item $A\cong A\otimes\Z$;
\item $\nucdim(A)\le1$.
\end{enumerate}
\end{theorem}

The implication from finite nuclear dimension to absorption proves
Conjecture (C1) of Robert and Tikuisis \cite{RT}. Together the two
implications establish the nuclear-dimension equivalence in the
nonsimple Toms--Winter regularity question of
Schafhauser--Tikuisis--White
\cite[Section~22, Problem~LXXV]{STW99}.
Here and throughout, unadorned tensor products are spatial. The
maximal tensor products used in the independent splitting theorem
below carry an explicit subscript.

The companion \cite[Corollaries~1.3 and~1.4]{CuntzCompanion} adds
Cuntz-semigroup criteria to the main equivalence.

\begin{corollary}[Regularity and Cuntz-semigroup criteria]
\label{cor:regularity-cu}
Under the hypotheses of Theorem~\ref{thm:main}, let
$\iota_A:A\to A\otimes\Z$ be $\iota_A(a)=a\otimes1_\Z$.
The conditions of that theorem are also equivalent to each of the following:
\begin{enumerate}[label=\textup{(\roman*)}]
\item the canonical map $\Cu(\iota_A):\Cu(A)\to\Cu(A\otimes\Z)$ is
an isomorphism in $\mathbf{Cu}$;
\item there is an abstract isomorphism $\Cu(A)\cong\Cu(A\otimes\Z)$
in $\mathbf{Cu}$.
\end{enumerate}
The abstract isomorphism need not be induced by a $*$-homomorphism.
\end{corollary}

The proof, including the nonunital naturality argument for the canonical
map, follows the proof of Theorem~\ref{thm:main} in
Section~\ref{subsec:cu-consequence}.

The dimension construction has a broader formulation, without a
separate assumption on elementary subquotients.

\begin{theorem}[Sharp bound after Jiang--Su absorption]
\label{intro:main}
For every separable nuclear complex $C^*$-algebra $A_0$,
\[
 \nucdim(A_0\otimes\Z)\le1.
\]
\end{theorem}

The bound is optimal uniformly over this class. Nuclear dimension zero
for a separable algebra implies the AF property
\cite[Remarks~2.2(iii)]{WZNuclear}, whereas $\Z$ is unital and
infinite-dimensional with no projections other than $0$ and $1$.
Thus $\nucdim(\Z)=1$.

The absorption implication uses a separate theorem. A positive element
is \emph{full} if it generates the whole ambient algebra as a closed
two-sided ideal, and a \emph{character} is a nonzero $*$-homomorphism
to $\C$.

\begin{theorem}[Full signs in a finite maximal tensor power]
\label{thm:tensor-intro}
For every nonzero unital complex $C^*$-algebra $D$ with no characters,
there is a finite integer $n\ge1$ such that
$D^{\otimes_{\max}n}$ contains two full orthogonal positive elements.
\end{theorem}

No separability or nuclearity is assumed in this theorem. Its finite
conclusion answers affirmatively Kirchberg and R\o rdam's
Question~6.2 \cite{KR}, which asks for this splitting property in the
infinite maximal tensor power. The passage to the infinite power is
by its canonical unital inclusion, which preserves fullness and
orthogonality.

The companion \cite[Theorems~1.1 and~1.2]{Fixed324} obtains another
finite endpoint under the same hypotheses on $D$: for some $m\ge1$,
there is a unital $*$-homomorphism from the prime dimension-drop algebra
$I(2,3)$ to $D^{\otimes_{\max}m}$. For every free ultrafilter $\omega$ on
$\N$, it also proves that a nonzero unital separable complex $C^*$-algebra
$A$ is $\Z$-stable exactly when its norm central-sequence algebra
$A_\omega\cap A'$ has no characters. The proof below uses only the full-sign
endpoint of Theorem~\ref{thm:tensor-intro}.

\subsection{Background and the two obstacles}

Winter and Zacharias introduced nuclear dimension as a noncommutative
covering dimension and established its basic permanence properties
\cite{WZNuclear}. Their order-zero structure theorem
\cite{WZOrderZero} makes a colour into a homomorphism multiplied by a
commuting positive weight. Toms and Winter's theory of strongly
self-absorbing algebras \cite{TW} explains why absorption can supply
commuting tensor factors without changing the isomorphism class.
In the simple setting, Matui--Sato proved that a unital separable simple
nuclear quasidiagonal algebra with a unique trace has decomposition rank
at most one after tensoring with the universal UHF algebra
\cite[Theorem~4.2]{MS}. Decomposition rank requires the same order-zero
decomposition as nuclear dimension, together with contractivity of the
entire return map; its bounds therefore also bound nuclear dimension.
Sato--White--Winter established the bound three
for separable unital simple nuclear $\Z$-stable algebras with a unique trace
\cite[Theorem~B]{SWW}. The two-coloured classification work of
Bosa--Brown--Sato--Tikuisis--White--Winter obtained nuclear dimension
one for non-AF separable unital simple nuclear $\Z$-stable algebras
whose trace simplex is nonempty and Bauer
\cite[Theorem~F]{BBSTWW}. Castillejos--Evington--Tikuisis--White--Winter
removed the trace-boundary restriction in the unital case
\cite[Theorem~B]{CETWW}, and Castillejos--Evington treated the
nonunital case \cite[Theorem~A]{CE}.

Nonsimple antecedents reduce dimension in specific structural
classes. Tikuisis--Winter proved decomposition rank at most two after
$\Z$-stabilization for algebras locally approximable by hereditary
subalgebras of $C(X,\mathcal K)$, with $X$ compact Hausdorff
\cite[Theorem~4.1]{TiW}.
Elliott--Niu--Santiago--Tikuisis extended the bound two to locally
subhomogeneous $\Z$-stable algebras \cite[Theorem~A]{ENST}.
In the purely infinite direction, Szab\'o proved
$\nucdim(A\otimes\mathcal O_\infty)\le3$ for every separable nuclear
$A$ \cite[Theorem~2.3]{SzaboOInfinity}. Bosa--Gabe--Sims--White
then established nuclear dimension at most one for separable nuclear
$\mathcal O_\infty$-stable algebras \cite[Theorem~A]{BGSW}.
Their homomorphism theorem gives the same bound for a nuclear
$\mathcal O_\infty$-stable homomorphism from a separable exact domain
to an arbitrary target \cite[Theorem~C]{BGSW}.

The nonsimple dimension problem has an additional difficulty: a trace
may be bounded on one ideal and infinite outside it, and stably finite
and purely infinite ideal quotients may coexist. The appropriate trace
cone therefore consists of all lower semicontinuous extended traces,
as developed by Elliott--Robert--Santiago \cite{ERS}. Our estimates are
relative to the support rank of a spectral cut generating the ideal
being tested. This rank measures the trace of the support, rather than
the trace value of the cut itself. When it is finite, a relative estimate
controls the error; rank zero forces the error rank to vanish. When it
is infinite, a separate ideal-containment argument forces the comparison
target to have infinite rank. The construction begins with the convex order-zero approximations of
Hirshberg--Kirchberg--White \cite{HKW}, in the refinement of
Brown--Carri\'on--White with asymptotically order-zero incoming maps
\cite{BCW}. It localizes them inside hereditary
corners, and rounds their convex weights to obtain one matrix-valued
order-zero return map.

The comparison stage uses two distinct forms of regularity. At each
coordinate, exactness and $\Z$-stability turn a uniform extended-rank
gap into Cuntz comparison, using R\o rdam's almost-unperforation theorem
\cite{Rordam} and Elliott--Robert--Santiago's comparison criterion
\cite[Proposition~6.2]{ERS}. In the supported relative commutant of the ultrapower,
we instead use abstract dimension functions. Gabe's order-zero
formulation \cite{Gabe} of the Kirchberg--R\o rdam cone criterion
\cite{KR05} supplies finite compression
sums from pointwise ideal containment. Packing these sums yields comparisons
inside that commutant. Fullness, divisibility in a full corner, and
almost unperforation verify a further hereditary-subalgebra condition.
The unitary-equivalence criterion of
Bosa--Brown--Sato--Tikuisis--White--Winter \cite{BBSTWW}, which adapts
Robert--Santiago's positive-element classification method \cite{RS},
turns the comparisons into two weighted conjugacies.

For the reverse implication, Winter proved that a separable simple
unital nonelementary algebra of finite nuclear dimension absorbs $\Z$
\cite[Corollary~6.3]{WinterAbsorption}. Tikuisis extended this to
separable algebras of finite nuclear dimension with finitely many ideals
and no elementary quotient of an ideal
\cite[Corollary~8.7]{TikuisisFiniteIdeals}. Robert--Tikuisis obtained
further nonsimple cases under restrictions on the primitive ideal space
and subquotients \cite[Corollaries~7.11 and~7.16]{RT}, and reduced absorption
for separable algebras of finite nuclear dimension to the existence of
two full orthogonal positive elements in the central sequence algebra
\cite[Theorem~1.2]{RT}. Winter--Zacharias's finite covering-number bound
\cite[Proposition~4.3]{WZNuclear} rules out characters in that algebra
under the nowhere-scattered hypothesis. The remaining step is precisely
a full splitting construction. Kirchberg and R\o rdam isolated the
character and tensor-splitting questions in \cite{KR}. Our finite tensor
construction supplies the needed splitting by combining a uniform
Gaussian sign estimate with pure-state excision. The latter is the
method of Akemann--Anderson--Pedersen \cite{AAP}; we include its short
proof using Kadison transitivity \cite{Kadison}.

Recent Cuntz-semigroup approaches provide a complementary perspective.
Antoine--Perera--Thiel--Vilalta relate finite nuclear dimension and the
Global Glimm Property to comparison and divisibility
\cite[Theorem~6.5]{APTV}; Perera--Thiel--Vilalta characterize absorption
through the corresponding regularity of central sequence algebras
\cite[Theorem~3.6]{PTV}. Our reverse implication instead uses the
finite covering bound and the full-sign theorem to meet the
Robert--Tikuisis criterion directly.

\subsection{Proof overview}

The two implications are independent until the final assembly. Figure~\ref{fig:proof}
records their interfaces.

For the dimension bound, we first work with a unital separable nuclear
$A$, without absorption. In Section~\ref{sec:model}, quasicentral ideal
cuts give finitely many families of hereditary pieces. The pieces are
orthogonal within each fixed threshold family, and averaging the
family compressions retains the prescribed tests. We fix the pieces
and their comparison bounds by finite sums of ideal generators before
choosing inner approximation errors. This ordering turns norm errors into
bounds relative to the support ranks of those generators, uniformly
over all traces for which the ranks are finite.

Section~\ref{sec:traces} places the matrix model in
$A\otimes\Z_1\otimes\Z_2\otimes\Z_3$. The first $\Z$ factor carries
the model. In the norm ultrapower, a positive contraction $q$ identifies
a part on which the model agrees with the original inclusion. The
second factor supplies arbitrarily many orthogonal copies of a norm-one
positive contraction. Relative trace estimates compare the error part
with this common part, and a spare spectral cut bounds the comparison
witnesses before they pass to the ultrapower.

Section~\ref{sec:commutant} constructs a supporting map and uses
approximate domination to pack the error into the common part. This
controls dimension functions in a supported relative commutant without
assuming that commutant exact. A full corner, its divisibility, and
almost unperforation give the hereditary-subalgebra condition needed
for unitary equivalence.
The constant third factor supplies a full-spectrum positive contraction
$k$. Functional-calculus comparison gives unitary conjugacies for the two
weights $k$ and $1-k$. Their sum gives exactly two colours for a corner
inclusion. In Section~\ref{sec:endgame}, stabilization and hereditary
permanence preserve those two colours, and unitization handles the
nonunital algebra. This proves Theorem~\ref{intro:main}.

For absorption, Section~\ref{sec:gaussian} constructs a real tensor
that takes both signs on every product of sufficiently separated pairs
of unit vectors. A Hamming packing provides exponentially many weakly
correlated Gaussian evaluations. Their failure probability decays fast
enough to control a net of all admissible configurations. In
Section~\ref{sec:tensor}, a finite commutator witness converts the
absence of characters into uniform separation of pure-state expectation
vectors. Excision transfers the tensor evaluations to every irreducible
representation of the maximal tensor product. The positive and negative
parts of the resulting selfadjoint element are therefore both full.

Finally, Section~\ref{sec:central} uses the norm central sequence algebra
$F(A)$, with its annihilator quotient in the nonunital case. Finite
covering excludes characters; a separable characterless subalgebra
containing $1_{F(A)}$ has a finite tensor splitting by
Theorem~\ref{thm:tensor-intro}. Commuting
reindexed copies give a unital homomorphism from that tensor power into
$F(A)$. Fullness survives this map, so Robert--Tikuisis's criterion
gives absorption and completes Theorem~\ref{thm:main}.

\begin{figure}[tbp]
\centering
\small
\begin{tabular}{@{}p{0.46\linewidth}@{\hspace{0.06\linewidth}}p{0.46\linewidth}@{}}
\centering \textbf{Dimension bound} &
\centering \textbf{Absorption}\tabularnewline[6pt]
\centering $A$ unital, separable, nuclear\newline
(later $A=\widetilde{A_0}\otimes\Z$) &
\centering $A$ separable, nowhere scattered,\newline $\nucdim(A)<\infty$\tabularnewline[4pt]
\centering $\Downarrow$ & \centering $\Downarrow$\tabularnewline[4pt]
\centering Localized matrix model;\newline
relative ranks and bounded witnesses &
\centering Finite covering makes $F(A)$ characterless;\newline
choose separable characterless $D\subseteq F(A)$\newline
with $1_D=1_{F(A)}$\tabularnewline[4pt]
\centering $\Downarrow$ & \centering $\Downarrow$\tabularnewline[4pt]
\centering Supported commutant comparison;\newline
hereditary condition and weights $k$, $1-k$ &
\centering Independent splitting theorem:\newline
every nonzero unital characterless $D$ has\newline
full orthogonal signs in some $D^{\otimes_{\max}n}$\tabularnewline[4pt]
\centering $\Downarrow$ & \centering $\Downarrow$\tabularnewline[4pt]
\centering Two-colour corner approximation;\newline
stabilization and hereditary permanence &
\centering Reindex commuting copies to obtain\newline
$D^{\otimes_{\max}n}\longrightarrow F(A)$ unital;\newline
the full signs pass to $F(A)$\tabularnewline[4pt]
\centering $\Downarrow$ & \centering $\Downarrow$\tabularnewline[4pt]
\centering $\nucdim(A_0\otimes\Z)\le1$ &
\centering Robert--Tikuisis criterion:\newline
$A\cong A\otimes\Z$\tabularnewline
\end{tabular}
\caption{The two proof branches. The left branch uses spatial tensor
products, first constructs the unital model, and then passes to
$A_0\otimes\Z$ by hereditary permanence. The splitting theorem in the
right column requires neither separability nor nuclearity of $D$;
its application to $F(A)$ uses the displayed unital map.
The branches meet only in the main equivalence.}
\label{fig:proof}
\end{figure}


\newpage
\section{Preliminaries}
\label{sec:preliminaries}

We collect the approximation, comparison and diagonalization tools for
the dimension bound, including comparison from a uniform rank gap for
all lower semicontinuous extended traces. The central-sequence inputs
for the independent absorption implication are stated in
Section~\ref{sec:central}.

All $C^*$-algebras are complex. An unadorned tensor product is spatial;
maximal tensor products are explicitly marked $\otimes_{\max}$.
Representations are nondegenerate.  We write $D_+$ for the positive cone of
a $C^*$-algebra $D$, $\widetilde D$ for its unitization, and
\[
 \her_D(h)=\overline{hDh}\quad(h\in D_+).
\]
For any subset $S\subset D$, let $\Ideal_D(S)$ denote the closed two-sided
ideal generated by $S$; for $x\in D$ we abbreviate
$\Ideal_D(x)=\Ideal_D(\{x\})$.  The ambient algebra is omitted when it
is clear.  For $\varepsilon>0$,
$(h-\varepsilon)_+$ denotes the usual spectral cut.  A full element is
one generating the whole ambient algebra as a closed ideal.  A nonzero
$h\in D_+$ is \emph{totally full} if $f(h)$ is full for every nonzero
$f\in C_0((0,\|h\|])_+$.  For a subset $S$ of an ambient algebra,
$S^\perp$ consists of the elements annihilating $S$ on both sides.

\subsection{Order zero and finite-dimensional approximation}

A completely positive contractive map is abbreviated to c.p.c.  A
completely positive map $\sigma:D\to E$ has order zero if
$\sigma(a)\sigma(b)=0$ whenever $a,b\in D_+$ and $ab=0$.
For a c.p.c. order-zero map $\sigma$, the structure theorem of Winter--Zacharias
\cite[Theorem~3.3 and Corollaries~4.1--4.2]{WZOrderZero} gives a
$*$-homomorphism $\sigma_{\mathrm{hom}}$ and a commuting positive
contraction $h$ in the multiplier algebra of $C^*(\sigma(D))$ such that
\[
 \sigma(a)=h\sigma_{\mathrm{hom}}(a).
\]
For $f\in C_0((0,1])_+$, the formula
$f(\sigma)(a)=f(h)\sigma_{\mathrm{hom}}(a)$ defines an order-zero map
into $E$, contractive when $\|f\|\le1$.  We write
$\sigma^{1/l}$ for $f(t)=t^{1/l}$.  In particular,
\begin{equation}
 \label{prelim:order-zero}
 \|\sigma(a)\sigma(b)\|\le\|ab\|\quad(a,b\in D),
 \qquad
 \sigma(a)\sigma(b)=\sigma(1)\sigma(ab)
 \quad\text{if $D$ is unital}.
\end{equation}
The first inequality follows by writing the product as
$h^2\sigma_{\mathrm{hom}}(ab)$.  Multiplication of a c.p.c. order-zero map by
a positive contraction commuting with its entire range again gives a
c.p.c. order-zero map.

\Needspace{13\baselineskip}
\begin{definition}[Winter--Zacharias {\cite[Definition~2.1]{WZNuclear}}]
\label{prelim:nuclear-dimension}
For an integer $d\ge0$, a $C^*$-algebra $D$ has nuclear dimension at
most $d$, written $\nucdim(D)\le d$, if for every finite
$\mathcal F\subset D$ and $\varepsilon>0$ there are finite-dimensional
algebras $E_0,\ldots,E_d$, a c.p.c. map
$\psi:D\to\bigoplus_{j=0}^dE_j$, and c.p.c. order-zero maps
$\varphi_j:E_j\to D$ such that
\[
 \left\|\sum_{j=0}^d\varphi_j\psi_j(a)-a\right\|<\varepsilon
 \qquad(a\in\mathcal F).
\]
The $d+1$ summands are called colours.  The sum of the return maps
need not be contractive; each coloured return map must be contractive.
\end{definition}

We use the following approximation theorem in hereditary subalgebras as
well as in unital algebras.

\begin{theorem}[Brown--Carri\'on--White]
\label{prelim:bcw}
For every separable nuclear $D$ there are finite-dimensional algebras
$F_n$ and c.p.c. maps
\[
 D\xrightarrow{\alpha_n}F_n\xrightarrow{\beta_n}D
\]
such that $\beta_n\alpha_n\to\id_D$ in point-norm,
$\|\alpha_n(a)\alpha_n(b)\|\to0$ for orthogonal $a,b\in D_+$, and
each $\beta_n$ is a finite convex combination of c.p.c. order-zero maps
$F_n\to D$.
\end{theorem}

The convex order-zero decomposition is due to
Hirshberg--Kirchberg--White \cite[Theorem~1.4]{HKW}; the asymptotic
orthogonality refinement used here is \cite[Theorem~3.1]{BCW}.
In particular, the map induced by the
$\alpha_n$ in $\prod_nF_n/\bigoplus_nF_n$ is order zero, so
\begin{equation}
 \label{prelim:asymptotic-product}
 \limsup_n\|\alpha_n(a)\alpha_n(b)\|\le\|ab\|
 \qquad(a,b\in D).
\end{equation}
Nuclearity passes to hereditary subalgebras and
to unitizations.  A c.p.c. map from a separable nuclear algebra to a
quotient has a c.p.c. lift by the Choi--Effros lifting theorem
\cite[Theorem~3.10]{ChoiEffros}.  Every completely positive map from a nuclear algebra is
nuclear, by composition with finite-dimensional approximations of its
domain.

\subsection{Comparison and extended traces}

Write $\mathcal K$ for the compact operators on a separable infinite-dimensional
Hilbert space. For $u,v\in(D\otimes\mathcal K)_+$, write $u\precsim v$ if
$s_n^*vs_n\to u$ in norm for some $s_n\in D\otimes\mathcal K$.
Mutual comparison is Cuntz equivalence.  The semigroup $\Cu(D)$ consists
of the resulting classes, with block-sum addition and the order induced
by $\precsim$.  Its subsemigroup $W(D)$ consists of classes represented
in finite matrices over $D$.  We use the following standard properties:
\[
 [u+v]\le[u]+[v],\qquad
 [u+v]=[u]+[v]\ \text{if }uv=0,\qquad
 [s^*s]=[ss^*].
\]
Positive domination implies Cuntz comparison; elements of $\her(v)_+$
are Cuntz below $v$.  If $u\in\Ideal(v)_+$, then, for every
$\varepsilon>0$, there is $N<\infty$ with
$[(u-\varepsilon)_+]\le N[v]$.  Also
\[
 [u]=\sup_{\varepsilon>0}[(u-\varepsilon)_+],
 \qquad [(u-\varepsilon)_+]\ll[u],
\]
where $x\ll y$ means that any increasing sequence with supremum at
least $y$ has a term dominating $x$.  Spectral cuts are cofinal below
$[u]$ for this relation.  Projection classes are compact, and if a
projection $p$ is Cuntz below $v$, then $\her(v)$ contains a projection
Murray--von Neumann equivalent to $p$ (with matrix amplification when
needed).  These properties hold without simplicity or unitality; see
\cite[Section~4]{Rordam} and \cite[Sections~2 and~4]{ERS}.

\begin{lemma}[Finite compressions of a full element]
\label{prelim:full-compressions}
If $v\in D_+$ is full and $D$ is unital, then there are
$y_1,\ldots,y_N\in D$ such that
\[
 \sum_{l=1}^N y_l^*vy_l=1_D.
\]
\end{lemma}

\begin{proof}
The ideal criterion above, applied to the projection $1_D$, gives
$[1_D]\le N[v]$ for some $N$.  A sufficiently accurate Cuntz
approximation, compressed to the $1_D$ corner, gives
$x_1,\ldots,x_N\in D$ with
$\|1_D-\sum_lx_l^*vx_l\|<1$.  The positive element
$z=\sum_lx_l^*vx_l$ is invertible.  Taking $y_l=x_lz^{-1/2}$
gives the identity.
\end{proof}

A dimension function on $W(D)$ is an additive order-preserving map
$d:W(D)\to[0,\infty]$ with $d(0)=0$; no continuity is implicit.
The semigroup is \emph{almost unperforated} when
$(m+1)x\le my$ implies $x\le y$, for $m\in\N$.

\begin{theorem}[R\o rdam's comparison criterion]
\label{prelim:comparison}
If $D\cong D\otimes\Z$, then $W(D)$ is almost unperforated.
More generally, if $W(D)$ is almost unperforated, $u,v\in D_+$,
$u\in\Ideal_D(v)$, and
\[
 d([u])<d([v])\quad\text{for every dimension function $d$ with }d([v])=1,
\]
then $u\precsim v$.
\end{theorem}

These are \cite[Theorem~4.5 and Corollary~4.7]{Rordam}.  The second
assertion uses abstract dimension functions, and requires neither
exactness nor a trace-space hypothesis.

A lower semicontinuous extended trace on $D$ is an additive homogeneous
map $\tau:D_+\to[0,\infty]$ satisfying
$\tau(x^*x)=\tau(xx^*)$ and norm lower semicontinuity.
We allow traces that are not densely finite.  Their extensions to finite
matrices use the unnormalized matrix trace, and their extensions to the
stabilization are understood.  For a positive contraction $z$, put
\[
 d_\tau(z)=\lim_{l\to\infty}\tau(z^{1/l}).
\]
Scaling a nonzero positive element to a contraction gives the same
definition for arbitrary $z\ge0$.  We use
$\tau(z)\le d_\tau(z)$ for contractions and
\begin{equation}
 \label{prelim:trace-cut}
 d_\tau((z-\eta)_+)\le\eta^{-1}\tau(z)\qquad(\eta>0).
\end{equation}

The next lemma follows from Elliott--Robert--Santiago's comparison
criterion \cite[Proposition~6.2]{ERS} in the presence of a uniform rank
gap.  If $\Cu(D)$ is almost unperforated, their criterion gives
$x\le y$ for $x,y\in\Cu(D)$ from the ideal containment
$x\le\sup_n ny$ and the strict inequalities
$\lambda(x)<\lambda(y)$ for every functional $\lambda$ with
$\lambda(y)=1$.  A functional is an additive order-preserving map to
$[0,\infty]$ that vanishes at zero and preserves increasing suprema.
For exact algebras, the extended trace correspondence converts the rank
gap into these inequalities, while the trace that is zero on an ideal
and infinite outside it detects ideal containment.  We give the direct
regularization proof.

\begin{lemma}[Comparison with a uniform rank gap]
\label{prelim:half-comparison}
Let $D$ be exact and $\Z$-stable.  If $u,v\in D_+$ satisfy
\[
 d_\tau(u)\le\tfrac12d_\tau(v)
 \quad\text{for every lower semicontinuous extended trace $\tau$ on $D$},
\]
then $u\precsim v$.
\end{lemma}

\begin{proof}
Let $I=\Ideal_D(v)$.  The map that is zero on $I_+$ and infinity on
$D_+\setminus I$ is a lower semicontinuous extended trace: additivity
follows from heredity of $I$, the trace identity from
$x^*x\in I\Longleftrightarrow xx^*\in I$, and lower semicontinuity
from closedness of $I$.  The assumed inequality for this trace forces
$u\in I$.

Work in $D\otimes\mathcal K$, whose $W$-semigroup identifies with
$\Cu(D)$ and is almost unperforated.  Fix a dimension function $d$ on
this semigroup with $d([v])=1$.  Its regularization is
\[
 \overline d([z])=\sup_{\delta>0}d([(z-\delta)_+]).
\]
By \cite[Lemma~4.7]{ERS} and the spectral-cut description of compact
containment, $\overline d$ is a functional on $\Cu(D)$.  The
correspondence in \cite[Proposition~4.2]{ERS} realizes it as
$d_\lambda$ for a lower semicontinuous extended $2$-quasitrace.
Haagerup proved that finite-valued $2$-quasitraces on exact unital
algebras are traces.  Blanchard--Kirchberg extended this conclusion to
lower semicontinuous $2$-quasitraces on exact algebras, allowing a
possibly proper finite-domain ideal \cite[Remarks~4.3 and~4.5]{ERS}.
Thus $\lambda$ is a trace.  Restrict $\lambda$ to the original corner
$D$.  For each fixed $\eta>0$ the hypothesis gives
\[
 \begin{split}
 d([(u-\eta)_+])
 &\le\overline d([u])=d_\lambda(u)
 \le\tfrac12d_\lambda(v)
 =\tfrac12\overline d([v])\\
 &\le\tfrac12<1=d([v]).
 \end{split}
\]
The strict inequality is against $d([v])$, so it remains valid when
$\overline d([v])=0$.  Theorem~\ref{prelim:comparison} now compares
every $(u-\eta)_+$ with $v$.  Taking the supremum over cuts gives
$u\precsim v$.  Finally, comparisons between elements of the original
corner can be witnessed in that corner, by compressing each witness on
both sides by its multiplier projection.
\end{proof}

\begin{lemma}[Extended traces on finite matrices]
\label{prelim:matrix-traces}
Let $A$ be unital, let $k\ge1$, and let $\lambda$ be a lower
semicontinuous extended trace on $A\otimes M_k$.  Set
$\tau(a)=\lambda(a\otimes1_k)$ and
$L_k=\id_A\otimes\operatorname{tr}_k$, with normalized matrix trace.
Then
\[
 \lambda(z)=\tau(L_k(z))\qquad(z\in(A\otimes M_k)_+).
\]
No finiteness assumption on $\lambda$ is needed.
\end{lemma}

\begin{proof}
Write $e_{ij}=1_A\otimes e_{ij}$.  Additivity and the trace identity give
\[
 \lambda(z)
 =\sum_{i=1}^k\lambda(z^{1/2}e_{ii}z^{1/2})
 =\sum_{i=1}^k\lambda(e_{ii}ze_{ii}).
\]
Unitary equivalence of the diagonal corners gives
$\lambda(a\otimes e_{ii})=k^{-1}\tau(a)$ for $a\in A_+$.
Summing proves the formula.  Only finite sums and multiplication by
positive scalars occur, so the argument also applies to infinite values.
\end{proof}

\subsection{Ultraproducts and bounded witnesses}

Fix a free ultrafilter $\omega$ on $\N$.  For a sequence $(D_n)$ we use
the norm ultraproduct
\[
 \prod_\omega D_n
 =\prod_nD_n\big/\{(x_n):\lim_{n\to\omega}\|x_n\|=0\}.
\]
For a constant sequence this is denoted $D_\omega$; the ordinary
sequence algebra is $D_\infty=\ell^\infty(D)/c_0(D)$.

The following form of Kirchberg's $\varepsilon$-test is recorded in
\cite[Lemma~1.10]{BBSTWW}; we include the diagonal proof.

\begin{lemma}[Kirchberg's $\varepsilon$-test]
\label{prelim:epsilon-test}
Let $X_n$ be nonempty sets and let
$f_{j,n}:X_n\to[0,\infty)$ be given for $j,n\in\N$.
Suppose that for every $m$ and $\delta>0$ there is
$(x_n)\in\prod_nX_n$ with
$\lim_{n\to\omega}f_{j,n}(x_n)<\delta$ for $1\le j\le m$.
Then there is a single $(x_n)$ for which all these ultralimits are zero.
\end{lemma}

\begin{proof}
Choose $(x_n^{(m)})$ solving the first $m$ conditions with ultralimits
less than $1/(2m)$.  Choose $A_m\in\omega$ on which the first $m$
errors of this choice are less than $1/m$, and set
$E_m=A_1\cap\cdots\cap A_m$.  These sets are decreasing and belong
to $\omega$.
At coordinate $n$, use $x_n^{(l(n))}$, where
$l(n)=\max\{m\le n:n\in E_m\}$, using any element of $X_n$ if the
set is empty.  Then $l(n)\to_\omega\infty$, and every fixed error tends
to zero.
\end{proof}

We apply this to countably many norm relations on contractions and on
variables with fixed, separately prescribed norm bounds.  Nuclear
c.p.c. lifting permits the same procedure for maps out of a separable
nuclear algebra; a countable dense set records linear relations,
complete positivity and the needed order-zero relations.  Countably
many extra commutators or norms after a fixed contractive map can be
included among the tests.  The following elementary estimate supplies
the bounds for Cuntz witnesses.

\begin{lemma}[A spare spectral cut gives bounded witnesses]
\label{prelim:bounded-witness}
Let $u,y\in D_+$ and $0\le a<b$.  If
$u\precsim(y-b)_+$, then for every $\varepsilon>0$ there is $r\in D$
such that
\[
 \|u-r^*(y-a)_+r\|<\varepsilon,
 \qquad
 \|r\|^2\le\frac{\|u\|+\varepsilon}{b-a}.
\]
Here $(y-0)_+=y$.  In particular the bound is uniform in a family of
contractions $u$ with a common gap $b-a>0$.
\end{lemma}

\begin{proof}
Choose $s$ with $\|u-s^*(y-b)_+s\|<\varepsilon$ and put
\[
 f(t)=\left(\frac{(t-b)_+}{(t-a)_+}\right)^{1/2},
 \qquad r=f(y)s,
\]
where the ratio is zero when its denominator vanishes.  This is a
continuous function vanishing near zero, and
\[
 f(y)(y-a)_+f(y)=(y-b)_+,
 \qquad f(t)^2\le(b-a)^{-1}(t-b)_+.
\]
Consequently
\[
 r^*(y-a)_+r=s^*(y-b)_+s,
 \qquad
 \|r\|^2\le(b-a)^{-1}\|s^*(y-b)_+s\|,
\]
which proves the claim.  If the original comparison is formulated in
the stabilization, compression to the original corner gives $s\in D$.
\end{proof}

The choices $(a,b)=(c,2c)$ and $(a,b)=(0,c)$ will both be used.
Together with Lemma~\ref{prelim:epsilon-test}, the estimate permits
coordinatewise comparisons with a fixed spare cut to pass to the norm
ultraproduct.  Coordinatewise Cuntz comparison without such control
does not itself supply bounded representatives.

\subsection{Domination and unitary equivalence}

\begin{theorem}[Gabe's approximate domination theorem]
\label{prelim:domination}
Let $D$ be exact and let $\Phi,\Psi:D\to E$ be nuclear completely
positive maps, with $\Psi$ of order zero.  Suppose
\[
 \Phi(a)\in\Ideal_E(\Psi(a))\qquad(a\in D_+).
\]
For every finite $\mathcal F\subset D$ and $\varepsilon>0$ there are
$y_1,\ldots,y_N\in E$ such that
\[
 \left\|\Phi(a)-\sum_{l=1}^Ny_l^*\Psi(a)y_l\right\|<\varepsilon
 \qquad(a\in\mathcal F).
\]
\end{theorem}

This is the implication from pointwise ideal containment to approximate
domination in \cite[Theorem~3.3]{Gabe}, building on the ideal-separation
method of \cite[Proposition~4.2]{KR05}.  The approximation is in norm
simultaneously on $\mathcal F$, and the codomain may be a nonexact norm
ultraproduct.  No bound on the individual $y_l$ is asserted.

\begin{lemma}[Relative-commutant unitary equivalence]
\label{prelim:unitary}
Let $D_n$ be separable unital $\Z$-stable algebras, put
$D_\omega=\prod_\omega D_n$, and let $A$ be separable and unital.
For a c.p.c. order-zero map $\rho:A\to D_\omega$ set
\[
 C=D_\omega\cap\rho(A)'\cap\{1-\rho(1)\}^{\perp}.
\]
Assume that $C$ is full in $D_\omega$ and that every full hereditary
subalgebra $E\subset C$ has the following property:
whenever $x\in E$ admits totally full $e_l,e_r\in E_+$ with
$e_lx=xe_r=0$, there is a full $s\in E$ satisfying $sx=xs=0$.
If $h_0,h_1\in C_+$ are totally full contractions and
\[
 f(h_0)\sim f(h_1)\quad\text{in }C
 \qquad\text{for every }f\in C_0((0,1])_+,
\]
then $uh_1u^*=h_0$ for some unitary $u\in\widetilde C$.
\end{lemma}

This is \cite[Lemma~5.3]{BBSTWW}, whose proof adapts the
positive-element classification method of Robert--Santiago \cite{RS};
$\sim$ here denotes Cuntz equivalence.  We will verify the hereditary-subalgebra condition
directly.  After adjoining the ambient unit, the implementing unitary
belongs to $D_\omega$ and commutes with $\rho(A)$.  More explicitly,
if $C$ already has a unit $p_C$ different from the ambient unit, replace
a unitary $u\in C$ by $u+(1-p_C)$; otherwise use the canonical map
from the unitization to $C+\C1_{D_\omega}$.

\section{Matrix models with local trace control}
\label{sec:model}
\label{model:section}

Throughout this section, $A$ is a nonzero separable unital nuclear
$C^*$-algebra; no absorption hypothesis is imposed.  We construct a
matrix model with trace estimates localized to finitely many ideals.  The point is that a norm error
inside a suitable hereditary corner becomes a trace error relative
to the rank of the positive element generating the relevant ideal.
For $k\geq 1$, write
\[
 L_k=\id_A\otimes\mathrm{tr}_k\colon A\otimes M_k\longrightarrow A,
 \qquad
 L_k((z_{rs})_{r,s=1}^k)=\frac1k\sum_{r=1}^k z_{rr},
\]
where $\mathrm{tr}_k$ is the normalized matrix trace.  We may choose
all matrix sizes to be powers of $2$; this convenient restriction
costs nothing in the finite probability-rounding argument below.

Each test uses a positive contraction $a$ and a number $0<\gamma<1$.
We keep the uncut element $a$ for the later comparison target, estimate
the model error on the deeper cut $x=(a-\gamma)_+$, and use the
shallower cut $b=(a-\gamma/2)_+$ to specify both the ideal and the
rank scale for that estimate.  The two cuts have distinct roles:
$x$ lies in the hereditary subalgebra generated by $b$, whose rank
$d_\tau(b)$ will control the trace error.

Choose a countable dense subset $\mathcal D$ of the positive unit ball
of $A$ which is closed under cuts $a\mapsto(a-r)_+$ for positive
rational $r$.  Such a set is obtained from any countable dense subset
by adjoining its successive rational cuts.  Enumerate all pairs
$(a,\gamma)\in\mathcal D\times((0,1)\cap\mathbb Q)$ as
$(a_j,\gamma_j)_{j\geq1}$, and set
\begin{equation}
 \label{model:tests}
 x_j=(a_j-\gamma_j)_+,
 \qquad b_j=(a_j-\gamma_j/2)_+,
 \qquad I_j=\Ideal(b_j).
\end{equation}
In particular, $x_j\in\her(b_j)$ and $d_\tau(x_j)\leq d_\tau(b_j)$
for every lower semicontinuous extended trace $\tau$.

\begin{lemma}[Matrix model]
 \label{model:construction}
There are finite-dimensional $C^*$-algebras $F_n$, powers $k_n$ of
$2$, and completely positive contractive maps
\[
 A\xrightarrow{\alpha_n}F_n
 \xrightarrow{\beta_n}A\otimes M_{k_n},
 \qquad \theta_n=\beta_n\alpha_n,
\]
with the following properties.
\begin{enumerate}
 \item Each $\beta_n$ has order zero, and
 \[
  \|\alpha_n(c)\alpha_n(d)\|\longrightarrow0
  \quad(c,d\in A_+,\ cd=0).
 \]
 \item For every $a\in A$,
 \[
  \|L_{k_n}(\theta_n(a))-a\|\longrightarrow0.
 \]
 \item For every $n,j\geq1$ there is a positive contraction
 $w_{n,j}\in A\otimes M_{k_n}$, and there are nonnegative numbers
 $\eta_{n,j},\delta_{n,j}$ tending to zero as $n\to\infty$ for each
 fixed $j$, such that
 \begin{align}
  \|w_{n,j}-\theta_n(x_j)\|&\leq\eta_{n,j},
       \label{model:retained-norm}\\
  \tau(L_{k_n}(w_{n,j}))
     &\leq \tau(x_j)+\delta_{n,j}d_\tau(b_j)
       \label{model:retained-trace}
 \end{align}
 for every lower semicontinuous extended trace $\tau$ on $A$ with
 $d_\tau(b_j)<\infty$.
\end{enumerate}
The coefficients in \eqref{model:retained-trace} are independent of
$\tau$.  One may take $\delta_{n,j}=2^{-n}$ for all $n,j$.
\end{lemma}

We first establish the finite-stage construction.  Its trace estimates
will be obtained within hereditary subalgebras on which the relevant
extended traces are bounded.

\begin{lemma}[Trace bounds in a hereditary subalgebra]
 \label{model:corner-traces}
Let $P,b\in A_+$ be contractions, and suppose that
$[P]\leq K[b]$ for a finite positive integer $K$.  If $\tau$ is a
lower semicontinuous extended trace with $d_\tau(b)<\infty$, then
its restriction to $D=\her(P)$ extends to a bounded positive linear
functional satisfying
\begin{equation}
 \label{model:corner-bound}
 \|\tau|_D\|\leq d_\tau(P)\leq Kd_\tau(b).
\end{equation}
Consequently, for $y,z\in D_+$,
\begin{equation}
 \label{model:corner-error}
 |\tau(y)-\tau(z)|
    \leq \|y-z\|K d_\tau(b).
\end{equation}
\end{lemma}

\begin{proof}
Choose an approximate unit $(e_l)$ for $D$ consisting of positive
contractive continuous functions of $P$ which vanish at zero.  Such
elements satisfy $d_\tau(e_l)\leq d_\tau(P)$.  For any positive
contraction $z\in D$, the elements $e_lze_l$ converge to $z$ in norm,
and
\[
 \tau(e_lze_l)\leq\tau(e_l^2)\leq d_\tau(P).
\]
Lower semicontinuity therefore gives $\tau(z)\leq d_\tau(P)$.
Monotonicity and additivity of $d_\tau$ on Cuntz classes give the
second inequality in \eqref{model:corner-bound}.  Thus $\tau$ is
bounded on the positive unit ball of $D$ and extends to the asserted
bounded positive functional.  Its norm bound proves
\eqref{model:corner-error}.  This argument also applies when
$d_\tau(b)=0$, in which case the restriction is identically zero.
\end{proof}

\begin{lemma}[Finite ideal slicing]
 \label{model:slicing}
Fix $m,N\geq1$, and let
$\mathcal T=\{1,\ldots,N\}^m$.  For every $t\in\mathcal T$ and
$S\subseteq\{1,\ldots,m\}$, there are positive contractions
$(P_{S,t,i})_{i\geq1}$ in $A$ with the following properties.
\begin{enumerate}
 \item For fixed $t,i$, the family $(P_{S,t,i})_S$ is orthogonal;
 in particular, $\sum_S P_{S,t,i}\leq1$.
 \item For every $a\in A$,
 $\|[P_{S,t,i},a]\|\to0$.  If $1\leq j\leq m$ and $j\notin S$, then
 \[
  \|P_{S,t,i}^{1/2}x_jP_{S,t,i}^{1/2}\|\longrightarrow0.
 \]
 \item If $j\in S$, then $P_{S,t,i}\in I_j$, and there is a finite
 positive integer $K_{S,t,i,j}$ such that
 \begin{equation}
  \label{model:finite-rank-bound}
  [P_{S,t,i}]\leq K_{S,t,i,j}[b_j].
 \end{equation}
 \item For every $a\in A$,
 \begin{equation}
  \label{model:slicing-approximation}
  \limsup_{i\to\infty}
  \left\|\frac1{N^m}\sum_{t\in\mathcal T}\sum_S
     P_{S,t,i}^{1/2}aP_{S,t,i}^{1/2}-a\right\|
  \leq\frac mN\|a\|.
 \end{equation}
\end{enumerate}
\end{lemma}

\begin{proof}
Write $A_\infty=\ell^\infty(A)/c_0(A)$, identifying $A$ with its
constant sequences.  There are commuting positive contractions
$u_1,\ldots,u_m\in A_\infty\cap A'$ such that $u_j$ is represented
by positive contractions in $I_j$ and $u_jx_j=x_j$.  Indeed, at
coordinate $i$ choose successively an element of a quasicentral
approximate unit of each $I_j$.  Require its commutators with the
first $i$ elements of a fixed dense subset of $A$ to have norm less
than $1/i$, require that it act on $x_j$ within $1/i$, and require
that it commute within $1/i$ with the elements already chosen at
that coordinate.  These are finitely many approximate-unit and
quasicentrality requirements.  They give exactly the stated
relations on passing to $A_\infty$.

Choose $N$ pairwise disjoint closed intervals contained in $(0,1)$.
For the $r$th interval choose continuous functions
$g_r^-,g_r^+\colon[0,1]\to[0,1]$ satisfying
\[
 g_r^-g_r^+=0,\qquad
 0\leq g_r^-+g_r^+\leq1,
\]
whose sum is one outside that interval.  We choose $g_r^-$ equal
to one below the interval and zero above it, and $g_r^+$ equal to
zero below the interval and one above it.  Within the interval the
two functions can decrease and increase, respectively, meeting at
zero.  Set $h_r=g_r^-+g_r^+$.  At every point of $[0,1]$ at most
one interval can contribute a loss, so
\begin{equation}
 \label{model:threshold-average}
 1-\frac1N\leq\frac1N\sum_{r=1}^N h_r\leq1.
\end{equation}
For $t=(t_1,\ldots,t_m)\in\mathcal T$ define
\[
 P_{S,t}=
  \prod_{j\in S}g_{t_j}^+(u_j)
  \prod_{j\in\{1,\ldots,m\}\setminus S}g_{t_j}^-(u_j).
\]
These are central positive contractions, and for fixed $t$ they
are pairwise orthogonal as $S$ varies.  Families for different values
of $t$ may overlap and are combined by averaging.  Joint functional
calculus and independent
averaging of the coordinates of $t$ give
\begin{align}
 Q&=\frac1{N^m}\sum_{t\in\mathcal T}\sum_S P_{S,t}
   =\prod_{j=1}^m\left(\frac1N\sum_{r=1}^N h_r(u_j)\right),
       \label{model:mean-slice}\\
 (1-1/N)^m1&\leq Q\leq1,
 \qquad \|1-Q\|\leq m/N.
       \label{model:mean-slice-bound}
\end{align}
If $j\in S$, then $P_{S,t}$ belongs to the image of
$\ell^\infty(I_j)$ in $A_\infty$, since $g_{t_j}^+(0)=0$.
If $1\leq j\leq m$ and $j\notin S$, then
$g_{t_j}^-(u_j)x_j=g_{t_j}^-(1)x_j=0$, whence $P_{S,t}x_j=0$.

We construct the required orthogonal representatives directly,
before imposing exact ideal membership. Start with positive
contractive representatives of a finite orthogonal family, whose
pairwise products tend to zero.  To separate a first representative
$a_i$ from a later representative $b_i$, choose $r_i>0$ tending to
zero slowly enough that $\|a_ib_i\|/r_i\to0$.  Let
$f_i\colon[0,1]\to[0,1]$ be one on $[0,r_i/2]$ and zero on
$[r_i,1]$, with $1-f_i(t)\leq2t/r_i$.  Replace the first element
by $(a_i-r_i)_+$ and every later one by $f_i(a_i)b_if_i(a_i)$.
These replacements are orthogonal to the first cut, and
\[
 \|b_i-f_i(a_i)b_if_i(a_i)\|
    \leq4\|a_ib_i\|/r_i\longrightarrow0.
\]
Choose $r_i$ using the maximum of the finitely many relevant
products.  Repeating with the later elements preserves
orthogonality already obtained and changes no quotient element.
This gives, for every $t$, orthogonal positive contractive lifts
$R_{S,t,i}$ of $P_{S,t}$.

For $j\in S$, the distance from $R_{S,t,i}$ to $I_j$ tends to zero.
Let $\varepsilon_i>0$ tend to zero and dominate all these distances,
over the finite collection of triples $(S,t,j)$ under
consideration.  The first cuts
\[
 T_{S,t,i}=(R_{S,t,i}-2\varepsilon_i)_+
\]
belong exactly to every $I_j$ with $j\in S$: their images in $A/I_j$
vanish by functional calculus.  Choose another sequence
$\sigma_i>0$ tending to zero, and put
\[
 P_{S,t,i}=(T_{S,t,i}-\sigma_i)_+.
\]
Because $T_{S,t,i}\in\Ideal(b_j)$, the usual ideal criterion for
Cuntz comparison gives a finite integer $K_{S,t,i,j}$ for which
\eqref{model:finite-rank-bound} holds.  Here the strict positivity
of the second cut is essential: every positive cut of an element
in $\Ideal(b_j)$ is Cuntz below a finite direct sum of copies of
$b_j$.  Both cuts preserve orthogonality and change the
representatives by at most $2\varepsilon_i+\sigma_i\to0$.
Consequently they preserve all the quotient relations established
above, while imposing all the ideal and finite-multiple conditions
simultaneously.

Finally, centrality of $P_{S,t}$ gives
$P_{S,t}^{1/2}aP_{S,t}^{1/2}=P_{S,t}a$ in $A_\infty$.
Thus \eqref{model:mean-slice-bound} implies
\eqref{model:slicing-approximation}, and $P_{S,t}x_j=0$ proves the
required vanishing of the omitted compressions.  No bound on
$K_{S,t,i,j}$ uniform in $i$ has been asserted or used.
\end{proof}

\begin{lemma}[Finite-stage model]
 \label{model:finite-stage}
Fix $m\geq1$, a finite subset $\mathcal F$ of the unit ball of $A$,
a finite family $\mathcal O$ of pairs of orthogonal positive
contractions in $A$, and $\varepsilon>0$.  There are a
finite-dimensional $C^*$-algebra $F$, a power $k$ of $2$, completely
positive contractive maps
\[
 A\xrightarrow{\alpha}F\xrightarrow{\beta}A\otimes M_k,
 \qquad\theta=\beta\alpha,
\]
and positive contractions $w_1,\ldots,w_m\in A\otimes M_k$ such
that $\beta$ has order zero and
\begin{align}
 \|L_k(\theta(a))-a\|&<\varepsilon
       &&(a\in\mathcal F), \label{model:finite-approx}\\
 \|\alpha(c)\alpha(d)\|&<\varepsilon
       &&((c,d)\in\mathcal O), \label{model:finite-orth}\\
 \|w_j-\theta(x_j)\|&<\varepsilon
       &&(1\leq j\leq m), \label{model:finite-retained}\\
 \tau(L_k(w_j))&\leq\tau(x_j)+\varepsilon d_\tau(b_j)
       &&(1\leq j\leq m, d_\tau(b_j)<\infty).
       \label{model:finite-trace}
\end{align}
The last inequality holds for every lower semicontinuous extended
trace $\tau$ satisfying the indicated finite-rank condition.
\end{lemma}

\begin{proof}
\emph{Fixing the corners.}
Choose $N$ so large that $m/N<\varepsilon/4$, and use
Lemma~\ref{model:slicing}.  Write
\[
 l_{S,t,i}(a)=P_{S,t,i}^{1/2}aP_{S,t,i}^{1/2}.
\]
Asymptotic centrality implies, for $(c,d)\in\mathcal O$,
\[
 \|l_{S,t,i}(c)l_{S,t,i}(d)\|
 =\|P_{S,t,i}^{1/2}cP_{S,t,i}dP_{S,t,i}^{1/2}\|
 \longrightarrow0.
\]
Together with \eqref{model:slicing-approximation} and the omitted
compression estimates, this permits us to fix one coordinate $i$
such that
\begin{align}
 \left\|\frac1{N^m}\sum_{t,S}l_{S,t,i}(a)-a\right\|
   &<\varepsilon/3 &&(a\in\mathcal F),
       \label{model:fixed-average}\\
 \|l_{S,t,i}(x_j)\|&<\varepsilon
       &&(1\leq j\leq m,\ j\notin S), \label{model:fixed-omitted}\\
 \|l_{S,t,i}(c)l_{S,t,i}(d)\|&<\varepsilon/2
       &&((c,d)\in\mathcal O). \label{model:fixed-products}
\end{align}
All these requirements range over finite sets.  From now on this
coordinate is fixed.  For each occurrence indexed by $(S,t)$,
suppress $i$ and write
\[
 P=P_{S,t,i},\quad D_P=\her(P),\quad
 l_P(a)=P^{1/2}aP^{1/2},\quad K_{P,j}=K_{S,t,i,j}.
\]
Every subscript $P$ below retains this $(S,t)$ occurrence index.
Thus coincident positive elements, including zero pieces, still
receive separate corner approximations and outcome labels.
In particular, the finitely many constants $K_{P,j}$ have already
been fixed before any approximation accuracy within $D_P$ is
chosen.

\emph{Approximating within the corners.}
Each $D_P$ is separable and nuclear.  The Brown--Carri\'on--White
approximation theorem, Theorem~\ref{prelim:bcw}, supplies
finite-dimensional approximations with asymptotically order-zero
first maps and return maps which are finite convex combinations
of completely positive contractive order-zero maps.
For a fixed corner, take a sequence of these first maps
$\zeta_r\colon D_P\to E_r$. The asymptotic product estimate
\eqref{prelim:asymptotic-product} gives
\begin{equation}
 \label{model:quotient-product-bound}
 \limsup_{r\to\infty}\|\zeta_r(y)\zeta_r(z)\|
   \leq\|yz\|\qquad(y,z\in D_P).
\end{equation}
Thus sufficiently late first maps have product norm less than
$\varepsilon$ on the finitely many pairs
$(l_P(c),l_P(d))$ in \eqref{model:fixed-products}.  This uses the
approximation sequence for the fixed corner; it requires no
uniformity in a varying hereditary subalgebra.

Choose positive numbers $\xi_P$ such that
\begin{equation}
 \label{model:inner-tolerances}
 \max_P\xi_P<\varepsilon/3,
 \qquad
 \frac1{N^m}\sum_{t\in\mathcal T}\sum_{S\ni j}
       K_{P_{S,t,i},j}\xi_{P_{S,t,i}}<\varepsilon/2
       \quad(1\leq j\leq m).
\end{equation}
There are only finitely many constants in these requirements, so
such choices are possible.  Choose a corresponding approximation
in each $D_P$, denoted
\[
 D_P\xrightarrow{\zeta_P}E_P
   \xrightarrow{\varphi_P}D_P,
 \qquad
 \varphi_P=\sum_{r=1}^{h_P}\lambda_{P,r}\beta_{P,r},
\]
where the $\lambda_{P,r}$ are nonnegative and sum to one, and each
$\beta_{P,r}$ is completely positive contractive and order zero.
All component return maps for this corner have the same domain
$E_P$ and use the same first map $\zeta_P$.  Arrange that
\begin{align}
 \|\varphi_P\zeta_Pl_P(a)-l_P(a)\|&<\xi_P
       &&\bigl(a\in\mathcal F\cup\{x_1,\ldots,x_m\}\bigr),
       \label{model:inner-approx}\\
 \|\zeta_Pl_P(c)\zeta_Pl_P(d)\|&<\varepsilon
       &&((c,d)\in\mathcal O).
       \label{model:inner-products}
\end{align}
For a zero corner take $E_P=0$ and the zero maps, with a single
return component of weight one.

\emph{Individual outcomes and their retained parts.}
Consider the finite probability space whose outcomes $\nu$
consist of a tuple $t\in\mathcal T$ and, for each subset $S$, a
choice $r_S\in\{1,\ldots,h_{P_{S,t,i}}\}$.  Give $t$ uniform
probability and make the return-map choices independently with
the probabilities $\lambda_{P,r}$.  Write $\lambda_\nu$ for
the resulting probability, and let
\[
 E_t=\bigoplus_S E_{P_{S,t,i}},\qquad
 \alpha_t(a)=\bigoplus_S\zeta_{P_{S,t,i}}l_{P_{S,t,i}}(a),
\]
\[
 \beta_\nu((z_S)_S)=
   \sum_S\beta_{P_{S,t,i},r_S}(z_S),
 \qquad T_\nu=\beta_\nu\alpha_t.
\]
The map $\alpha_t$ is completely positive contractive because a
finite direct sum has the maximum norm.  The hereditary
subalgebras $D_{P_{S,t,i}}$ are mutually orthogonal for the single
tuple $t$ used by the outcome.
Consequently $\beta_\nu$ is completely positive contractive
and order zero: each summand has these properties, and products
between distinct summand ranges vanish.  Moreover,
\eqref{model:inner-products} bounds the product defect of every
$\alpha_t$ by $\varepsilon$, uniformly over the outcomes.

Independence of the return-map choices gives
\begin{equation}
 \label{model:expected-map}
 \sum_\nu\lambda_\nu T_\nu(a)
   =\frac1{N^m}\sum_{t\in\mathcal T}\sum_S
       \varphi_{P_{S,t,i}}\zeta_{P_{S,t,i}}l_{P_{S,t,i}}(a).
\end{equation}
For each fixed $t$, the differences in
\eqref{model:inner-approx} lie in orthogonal hereditary
subalgebras, so their sum has norm at most $\max_P\xi_P$.
Equations \eqref{model:fixed-average},
\eqref{model:inner-tolerances}, and \eqref{model:expected-map}
therefore imply
\begin{equation}
 \label{model:expected-approximation}
 \left\|\sum_\nu\lambda_\nu T_\nu(a)-a\right\|
    <2\varepsilon/3\qquad(a\in\mathcal F).
\end{equation}

For $1\leq j\leq m$, retain only the corners indexed by $S\ni j$:
\begin{equation}
 \label{model:outcome-retained}
 W_{\nu,j}=
    \sum_{S\ni j}\beta_{P_{S,t,i},r_S}
                \zeta_{P_{S,t,i}}l_{P_{S,t,i}}(x_j).
\end{equation}
This is a positive contraction.  Contractivity of each component
and orthogonality of the omitted corners give, for every
individual outcome,
\begin{equation}
 \label{model:before-averaging}
 \|T_\nu(x_j)-W_{\nu,j}\|
   \leq\max_{S\not\ni j}\|l_{P_{S,t,i}}(x_j)\|
   <\varepsilon.
\end{equation}
The maximum over an empty family is taken to be zero.  In
particular, the estimate holds before taking any average.

\emph{The uniform relative trace estimate.}
Fix $j$ and a lower semicontinuous extended trace $\tau$ with
$d_\tau(b_j)<\infty$.  In every retained corner,
Lemma~\ref{model:corner-traces} applies with the already fixed
constant $K_{P,j}$.  Thus \eqref{model:inner-approx} implies
\[
 \tau(\varphi_P\zeta_Pl_P(x_j))
    \leq\tau(l_P(x_j))+\xi_PK_{P,j}d_\tau(b_j)
    \qquad(j\in S).
\]
All quantities in this inequality are finite.  For each fixed
$t$, traciality and $\sum_{S\ni j}P_{S,t,i}\leq1$ give
\begin{align*}
 \sum_{S\ni j}\tau(l_{P_{S,t,i}}(x_j))
 &=\tau\left(x_j^{1/2}
       \Bigl(\sum_{S\ni j}P_{S,t,i}\Bigr)x_j^{1/2}\right)\\
 &\leq\tau(x_j).
\end{align*}
Taking expectations in \eqref{model:outcome-retained} and using
\eqref{model:inner-tolerances} proves
\begin{equation}
 \label{model:expected-trace}
 \sum_\nu\lambda_\nu\tau(W_{\nu,j})
    \leq\tau(x_j)+(\varepsilon/2)d_\tau(b_j).
\end{equation}
For later use there is also an individual bound
\begin{equation}
 \label{model:individual-trace}
 \tau(W_{\nu,j})\leq M_{\nu,j}d_\tau(b_j),
 \qquad M_{\nu,j}=\sum_{S\ni j}K_{P_{S,t,i},j}<\infty.
\end{equation}
Every coefficient here was chosen independently of $\tau$.
If $d_\tau(b_j)=0$, then each retained output has trace exactly
zero, as does $x_j$.  Hence the same inequalities hold in this
case without an absolute additive error.  No condition is being
placed on traces for which $d_\tau(b_j)=\infty$.

\emph{Rounding the probabilities and forming the matrix model.}
After all corner maps and the constants $M_{\nu,j}$ have been
fixed, choose $k=2^s$ and nonnegative integers $n_\nu$ with
$\sum_\nu n_\nu=k$ so that
$\mu_\nu=n_\nu/k$ satisfies
\begin{equation}
 \label{model:rounding-errors}
 \sum_\nu|\mu_\nu-\lambda_\nu|<\varepsilon/3,
 \qquad
 \sum_\nu|\mu_\nu-\lambda_\nu|M_{\nu,j}
      <\varepsilon/2\quad(1\leq j\leq m).
\end{equation}
Such choices exist for arbitrarily large $s$: round the finitely
many numbers $k\lambda_\nu$ down to integers and distribute
the remaining units among the outcomes.  The resulting
probabilities stay in the probability simplex, and their total
absolute error tends to zero as $s\to\infty$.  Contractivity of
the $T_\nu$ shows that replacing $\lambda_\nu$ by
$\mu_\nu$ changes the expected map on the unit ball by less
than $\varepsilon/3$.  Equations
\eqref{model:individual-trace} and \eqref{model:rounding-errors}
show that the expected retained trace changes in absolute value
by at most $(\varepsilon/2)d_\tau(b_j)$.

List the outcomes with their multiplicities as
$\nu_1,\ldots,\nu_k$, and write $t_l$ for the threshold
tuple in $\nu_l$.  Define
\[
 F=\bigoplus_{l=1}^k E_{t_l},\qquad
 \alpha(a)=\bigoplus_{l=1}^k\alpha_{t_l}(a),\qquad
 \beta((z_l)_{l=1}^k)
    =\operatorname{diag}\bigl(\beta_{\nu_l}(z_l)\bigr)_{l=1}^k.
\]
The map $\alpha$ is completely positive contractive, and $\beta$
is completely positive contractive and order zero, by the
orthogonality of the matrix diagonal positions.  We have
\[
 \theta(a)=\operatorname{diag}(T_{\nu_l}(a))_{l=1}^k,
 \qquad
 L_k(\theta(a))=\sum_\nu\mu_\nu T_\nu(a).
\]
Consequently \eqref{model:expected-approximation} and
\eqref{model:rounding-errors} give \eqref{model:finite-approx},
and the maximum norm in the direct sum gives
\eqref{model:finite-orth}.  Finally, put
\[
 w_j=\operatorname{diag}(W_{\nu_l,j})_{l=1}^k.
\]
These are positive contractions, and the estimate for individual
outcomes in \eqref{model:before-averaging} proves
\eqref{model:finite-retained}.  Since
$L_k(w_j)=\sum_\nu\mu_\nu W_{\nu,j}$,
\eqref{model:expected-trace} and the trace part of the rounding
estimate prove \eqref{model:finite-trace}.
\end{proof}

\begin{proof}[Proof of Lemma~\ref{model:construction}]
Fix a countable dense subset $(c_l)_{l\geq1}$ of the unit ball of
$A$.  The space
\[
 \{(c,d)\in A_+\times A_+:
           \|c\|,\|d\|\leq1,\ cd=0\}
\]
is a separable metric space, so choose a countable dense family
$(c'_l,d'_l)_{l\geq1}$ in this space.  At stage $n$, apply
Lemma~\ref{model:finite-stage} to the first $n$ tests,
$\mathcal F_n=\{c_1,\ldots,c_n\}$,
$\mathcal O_n=\{(c'_1,d'_1),\ldots,(c'_n,d'_n)\}$, and
$\varepsilon=2^{-n}$.  Denote the resulting maps by $\alpha_n$
and $\beta_n$, and the retained elements by $w_{n,j}$ for
$j\leq n$.

Contractivity extends \eqref{model:finite-approx} from the dense
family to point-norm convergence on $A$.  Likewise, if $c,d$ are
orthogonal positive contractions and $(c',d')$ is one of the
chosen orthogonal pairs, then
\[
 \|\alpha_n(c)\alpha_n(d)\|
 \leq\|c-c'\|+\|d-d'\|
       +\|\alpha_n(c')\alpha_n(d')\|.
\]
First taking the limit superior in $n$, then approximating
$(c,d)$ by such pairs, proves asymptotic preservation of
orthogonality.  Scaling proves the assertion for arbitrary
orthogonal positive elements.

For $j>n$ set $w_{n,j}=0$.  Take $\delta_{n,j}=2^{-n}$ for all
$n,j$, and set
\[
 \eta_{n,j}=
 \begin{cases}
  2^{-n},&j\leq n,\\
  1,&j>n.
 \end{cases}
\]
Equations \eqref{model:retained-norm} and
\eqref{model:retained-trace} now hold for every $n,j$, and both
error sequences tend to zero for each fixed $j$.  At each stage
the trace inequality holds simultaneously for all the indicated
extended traces.  Thus the diagonalization involves only the
countable lists of elements, tests, and orthogonal pairs; it
requires neither a countable family of traces nor a restriction
on the ideal lattice of $A$.
\end{proof}

\begin{remark}
 \label{model:composite-order-zero}
The composites $\theta_n$ also asymptotically preserve
orthogonality.  Indeed, the order-zero product estimate gives
\[
 \|\theta_n(c)\theta_n(d)\|
 \leq\|\alpha_n(c)\alpha_n(d)\|
 \qquad(c,d\in A_+).
\]
Consequently they induce a c.p.c. order-zero map into
$\prod_n(A\otimes M_{k_n})/\bigoplus_n(A\otimes M_{k_n})$,
and into the corresponding norm ultraproduct.  The same conclusion
holds after composing the $\beta_n$ with any c.p.c. order-zero maps
into a fixed target algebra: the displayed product estimate remains
valid for each composite return map.
\end{remark}

\section{The common part and localized trace comparison}
\label{sec:traces}

The matrix model will now be placed in a Jiang--Su tensor factor by an
order-zero map.  Its expectation is still asymptotically the identity.
This produces a common part of the model and the original inclusion.
The remaining part is small in a sense that can be compared uniformly
over all lower semicontinuous extended traces, including traces that are
infinite away from an ideal.

\subsection{Order-zero maps into the Jiang--Su algebra}

Write $\tau_{\Z}$ for the unique tracial state of $\Z$.
We first construct order-zero matrix maps and positive contractions
with arbitrarily many orthogonal copies.

\begin{lemma}\label{supplies:order-zero}
For every integer $k\geq 1$ and every $\xi>0$ there is a c.p.c.
order-zero map $\iota\colon M_k\to\Z$ such that
\[
 \tau_{\Z}(\iota(1))>1-\xi,
 \qquad \|\iota(e_{11})\|=1.
\]
\end{lemma}

\begin{proof}
Choose a multiple $d\geq2$ of $k$ with $(d+1)^{-1}<\xi$, and fix a
unital embedding $\mu\colon M_k\to M_d$.  Put
$p=e_{d+1,d+1}\in M_{d+1}$.  In $M_d\otimes M_{d+1}$, the complement
of $1\otimes p$ is the tensor product of two $d$-dimensional spaces.
Let $V$ be the flip on this complement, extended by the identity on
the range of $1\otimes p$.  The unitary group of each finite matrix
block is path connected.  There is therefore a continuous unitary path
$U(t)$, $0\leq t\leq1$, from $1$ to $V$, all of whose members commute
with $1\otimes p$.

Define
\[
 \iota_0(x)(t)
   =U(t)\bigl(\mu(x)\otimes(1-tp)\bigr)U(t)^*.
\]
This is a c.p.c. order-zero map into $C([0,1],M_d\otimes M_{d+1})$.
Its endpoint values are
\[
 \iota_0(x)(0)=\mu(x)\otimes1,
 \qquad
 \iota_0(x)(1)=1\otimes\operatorname{diag}(\mu(x),0).
\]
Consequently its range lies in the prime dimension-drop algebra
\[
 I_{d,d+1}
 =\{f\in C([0,1],M_d\otimes M_{d+1}):
 f(0)\in M_d\otimes1,\ f(1)\in1\otimes M_{d+1}\}.
\]
The commutation of $U(t)$ with $1\otimes p$ gives the exact defect
\[
 1-\iota_0(1)=\bigl(t\mapsto t(1\otimes p)\bigr).
\]
For $1\leq r\leq d+1$, the functions
\[
 y_r(t)=\sqrt t\,(1\otimes e_{r,d+1})
\]
belong to $I_{d,d+1}$ and satisfy
\[
 y_r^*y_r=1-\iota_0(1),
 \qquad
 \sum_{r=1}^{d+1}y_ry_r^*=(t\mapsto t1)\leq1.
\]
Thus every tracial state $\sigma$ on $I_{d,d+1}$ satisfies
\[
 (d+1)\sigma(1-\iota_0(1))
   =\sigma\left(\sum_{r=1}^{d+1}y_ry_r^*\right)\leq1.
\]
The prime dimension-drop algebra $I_{d,d+1}$ embeds
unitally into $\Z$ by \cite[Theorem~2.1(i)]{Rordam}.
Composing such an embedding with
$\iota_0$ gives the trace estimate.  Finally,
$\|\iota_0(e_{11})(0)\|=1$, so isometry of the embedding gives the
asserted norm.
\end{proof}

\begin{lemma}\label{supplies:copies}
For every integer $m\geq2$ there are a positive contraction
$D_m\in\Z$ of norm one, contractions $v_{m,1},\ldots,v_{m,m}\in\Z$,
and a finite positive integer $K_m$ such that
\begin{equation}\label{supplies:copy-relations}
 v_{m,r}^*v_{m,s}=\delta_{rs}D_m,
 \qquad
 1=\sum_{l=1}^{K_m}y_{m,l}^*(D_m-1/2)_+y_{m,l}
\end{equation}
for suitable $y_{m,l}\in\Z$.
\end{lemma}

\begin{proof}
Take $\iota\colon M_m\to\Z$ from
Lemma~\ref{supplies:order-zero}, and set
\[
 D_m=\iota(e_{11}),\qquad v_{m,r}=\iota^{1/2}(e_{r1}).
\]
Order-zero functional calculus and the product identity
\eqref{prelim:order-zero} give the first identity in
\eqref{supplies:copy-relations}.  The element $(D_m-1/2)_+$ is
nonzero, hence full in the simple unital algebra $\Z$.
Lemma~\ref{prelim:full-compressions} supplies the second identity.
\end{proof}

We fix these choices once and for all.  No uniform bound on $K_m$ is
required.

\subsection{The model and its common part}

Throughout this section, $A$ is a nonzero separable unital nuclear
$C^*$-algebra.  Set
\begin{equation}\label{setup:ambient}
 R=A\otimes\Z_1\otimes\Z_2\otimes\Z_3,
 \qquad B=R_\omega,
 \qquad B_1=(A\otimes\Z_1)_\omega\subset B.
\end{equation}
The first Jiang--Su factor carries the model and its common part;
the second supplies positive contractions with arbitrarily many
orthogonal copies. The constant third factor is reserved for the
complementary weights in the unitary-equivalence argument.
The tensor-factor inclusions are unital.  Write
$\pi\colon A\to B_1\subset B$ for the constant first-factor
inclusion. Nuclearity and strong self-absorption of $\Z$
\cite[Examples~1.14]{TW} imply that
$R$ is nuclear, exact, and $\Z$-stable.

Use the tests $a_j,\gamma_j,x_j,b_j$ from \eqref{model:tests}.
The matrix-model construction, Lemma~\ref{model:construction},
gives finite-dimensional algebras $F_n$, powers
$k_n$ of two, and c.p.c. maps
\[
 A\xrightarrow{\alpha_n}F_n
   \xrightarrow{\beta_n}A\otimes M_{k_n},
\]
where $\beta_n$ is order zero, $\alpha_n$ is asymptotically order zero,
and $L_{k_n}\beta_n\alpha_n(a)\to a$ for every $a\in A$.
Denote the lemma's retained positive contractions by
$W_{n,j}\in A\otimes M_{k_n}$. For each fixed $j$ they satisfy
\begin{equation}\label{setup:model-retained}
 \|W_{n,j}-\beta_n\alpha_n(x_j)\|\longrightarrow0
\end{equation}
and
\begin{equation}\label{setup:model-traces}
 \tau(L_{k_n}(W_{n,j}))
   \leq\tau(x_j)+\delta_{n,j}d_\tau(b_j)
 \qquad\bigl(d_\tau(b_j)<\infty\bigr).
\end{equation}
Here $\tau$ ranges over all lower semicontinuous extended traces on
$A$, and the nonnegative constants $\delta_{n,j}$ are independent of
$\tau$ and tend to zero for fixed $j$.

Choose maps $\iota_n\colon M_{k_n}\to\Z_1$ from
Lemma~\ref{supplies:order-zero} with
$s_n:=\tau_{\Z}(\iota_n(1))\to1$.  Define
\begin{equation}\label{setup:transformed-model}
 \begin{aligned}
 T_n&=\id_A\otimes\iota_n,&
 \beta'_n&=T_n\beta_n,\\
 \theta_n&=\beta'_n\alpha_n,&
 w_{n,j}&=T_n(W_{n,j}).
 \end{aligned}
\end{equation}
These maps and elements take values in $A\otimes\Z_1$, which we
regard as a subalgebra of $R$.

The order-zero product identity for $\iota_n$ tensors with
$\id_A$, first on elementary tensors and then by continuity.  Thus
$T_n$ and $\beta'_n$ are c.p.c. order-zero maps.  Put
\[
 L=\id_A\otimes\tau_{\Z}\colon A\otimes\Z_1\longrightarrow A.
\]
The functional $\tau_{\Z}\circ\iota_n$ is a positive trace on
$M_{k_n}$.  Indeed, the square-root order-zero map satisfies
\[
 \iota_n^{1/2}(z)^*\iota_n^{1/2}(z)=\iota_n(z^*z),
 \qquad
 \iota_n^{1/2}(z)\iota_n^{1/2}(z)^*=\iota_n(zz^*).
\]
Its mass is $s_n$, so
\[
 LT_n=s_nL_{k_n}.
\]
It follows that $L\theta_n(a)\to a$ for every $a\in A$.
The product estimate for an order-zero map also gives
\[
 \|\theta_n(c)\theta_n(d)\|
 \leq\|\alpha_n(c)\alpha_n(d)\|\longrightarrow0
 \qquad(c,d\in A_+,\ cd=0).
\]
Consequently the maps $\theta_n$ induce a c.p.c. order-zero map
\begin{equation}\label{setup:theta}
 \theta\colon A\longrightarrow B_1,
 \qquad \theta(a)=(\theta_n(a))_\omega,
 \qquad (w_{n,j})_\omega=\theta(x_j).
\end{equation}
The same letter $L$ will denote the induced unital completely positive
map $B_1\to A_\omega$.  With the constant inclusion of $A$ into
$A_\omega$ understood, we have $L\theta=L\pi=\id_A$.

\begin{lemma}[The common part]\label{setup:common}
There is a positive contraction $e\in B_1$ such that, with $q=1-e$,
\begin{equation}\label{setup:common-equations}
 L(e)=0,\qquad
 [e,\pi(a)]=[e,\theta(a)]=0,\qquad
 q\theta(a)=q\pi(a)
 \quad(a\in A).
\end{equation}
In particular, if $h=\theta(1)$, then $qh=q$.
\end{lemma}

\begin{proof}
Both ranges $\pi(A)$ and $\theta(A)$ belong to the multiplicative
domain of $L$.  For $\theta$, the two Schwarz inequalities give
\[
 a^*a=L(\theta(a))^*L(\theta(a))
 \leq L(\theta(a)^*\theta(a))
 \leq L(\theta(a^*a))=a^*a.
\]
The analogous equalities hold with $a$ and $a^*$ interchanged.
The assertion for $\pi$ follows in the same way, or directly from
its multiplicativity.

Let $D_0=C^*(\pi(A),\theta(A))\subset B_1$.  This is a separable
unital algebra, $L|_{D_0}$ is a homomorphism, and
\[
 J=\ker(L|_{D_0})
\]
is an ideal containing $\theta(a)-\pi(a)$ for every $a\in A$.
Choose a positive contractive quasicentral approximate unit $(f_l)$
for $J$ in $D_0$.  Then $L(f_l)=0$, the commutators of $f_l$ with
both ranges tend to zero, and
\[
 (1-f_l)(\theta(a)-\pi(a))\longrightarrow0
 \qquad(a\in A).
\]
Lift the $f_l$ to positive contractive sequences in
$A\otimes\Z_1$.  Apply the countable diagonalization of
Lemma~\ref{prelim:epsilon-test} to the preceding equations on dense
countable subsets, together with the coordinate test
$\|L(f_{l,n})\|$.  The latter tends to zero along $\omega$ for each
fixed $l$.  The resulting positive contraction $e\in B_1$ satisfies
\eqref{setup:common-equations}.  Taking $a=1$ gives $qh=q$.
\end{proof}

Fix positive contractive representatives
$e_n\in A\otimes\Z_1$ of $e$, and write
\begin{equation}\label{setup:representatives}
 q_n=1-e_n,\qquad \ell_n=\|L(e_n)\|.
\end{equation}
Thus $\ell_n\to_\omega0$.  Exact vanishing of $L(e_n)$ at any
coordinate is neither asserted nor needed.

\subsection{Rank estimates on a prescribed ideal}

Fix $j$ temporarily, and abbreviate
\[
 a=a_j,\quad \gamma=\gamma_j,\quad
 x=x_j,\quad b=b_j,\quad w_n=w_{n,j}.
\]
Inside $A\otimes\Z_1\subset R$, set
\begin{equation}\label{trace:coordinate-positives}
 \begin{split}
 u_{n,0}&=x^{1/2}e_nx^{1/2},\\
 u_{n,1}&=w_n^{1/2}e_nw_n^{1/2},\\
 z_n&=a^{1/2}(q_n-1/2)_+a^{1/2}.
 \end{split}
\end{equation}
These are positive contractions.

\begin{lemma}[Localized rank estimates]\label{trace:rank}
For each $j$ there are nonnegative numbers
$\varepsilon_{n,j}\to_\omega0$, independent of the trace, with the
following properties.  Let $\lambda$ be any lower semicontinuous
extended trace on $R$, and let $\tau=\lambda|_A$.  If
$d_\tau(b)<\infty$, then, for $i=0,1$ and $\eta>0$,
\begin{equation}\label{trace:small-rank}
 d_\lambda((u_{n,i}-\eta)_+)
   \leq\eta^{-1}\varepsilon_{n,j}d_\tau(b).
\end{equation}
If $\ell_n\leq1/4$, then, with no finiteness assumption,
\begin{equation}\label{trace:large-rank}
 d_\lambda((z_n-\gamma/16)_+)
   \geq(\gamma/16)d_\tau(b).
\end{equation}
\end{lemma}

\begin{proof}
We first work under the hypothesis $d_\tau(b)<\infty$.
On the hereditary subalgebra
$\her_A(b)\otimes\Z_1$ one has
\begin{equation}\label{trace:localized-expectation}
 \lambda(y)=\tau(L(y))\qquad(y\geq0).
\end{equation}
To justify this formula, every positive contraction in
$\her_R(b\otimes1)$ has $\lambda$-value at most
$d_\lambda(b\otimes1)=d_\tau(b)$.  The restriction of $\lambda$
to $\her_A(b)\otimes\Z_1$ therefore extends to a bounded trace;
the same holds for $\tau$ on $\her_A(b)$.  For a positive
$s\in\her_A(b)$, the functional
\[
 v\longmapsto\lambda(s\otimes v)
\]
is a bounded positive trace on $\Z_1$, of mass $\tau(s)$.
Its trace identity follows by applying the trace identity for
$\lambda$ to $s^{1/2}\otimes z$.  Uniqueness of the tracial state
of $\Z$ gives $\lambda(s\otimes v)=\tau(s)\tau_{\Z}(v)$.
Linearity and density prove \eqref{trace:localized-expectation}.
In particular,
\begin{equation}\label{trace:first-small-part}
 \lambda(u_{n,0})
 =\tau(x^{1/2}L(e_n)x^{1/2})
 \leq\ell_n\tau(x)
 \leq\ell_n d_\tau(b).
\end{equation}

There is a separate estimate for $w_n$, which need not lie in that
hereditary subalgebra.  Define
\[
 \rho_n=\lambda\circ T_n,
 \qquad
 \sigma_n(s)=\rho_n(s\otimes1)\quad(s\in A_+).
\]
Order-zero square-root calculus shows that $\rho_n$ is an extended
trace; its lower semicontinuity follows from continuity of $T_n$.
Furthermore,
\[
 \sigma_n(s)
 =\lambda(s\otimes\iota_n(1))\leq\lambda(s\otimes1)=\tau(s).
\]
Lemma~\ref{prelim:matrix-traces}, applied to $\rho_n$, gives
\begin{equation}\label{trace:matrix-formula}
 \rho_n(Y)=\sigma_n(L_{k_n}(Y))
 \qquad(Y\in(A\otimes M_{k_n})_+).
\end{equation}
This identity includes infinite values. It follows from
\eqref{setup:model-traces} that
\begin{equation}\label{trace:retained-finite}
 \lambda(w_n)
 \leq\tau(L_{k_n}(W_{n,j}))
 \leq\tau(x)+\delta_{n,j}d_\tau(b)<\infty.
\end{equation}

Put
\[
 S_n=w_n^{1/2}q_nw_n^{1/2},\qquad
 S'_n=x^{1/2}q_nx^{1/2},\qquad
 \Delta_{n,j}=\|S_n-S'_n\|.
\]
The common-part identity implies
$\Delta_{n,j}\to_\omega0$: in the quotient, the two expressions
are respectively $q\theta(x)$ and $q\pi(x)$.
Compressing $S'_n\leq S_n+\Delta_{n,j}1$ by $x^{1/(2r)}$ gives
\[
 \lambda(x^{1/(2r)}S'_nx^{1/(2r)})
 \leq\lambda(S_n)+\Delta_{n,j}\tau(x^{1/r})
 \leq\lambda(S_n)+\Delta_{n,j}d_\tau(b).
\]
Here the compression of $S_n$ has trace at most $\lambda(S_n)$,
by traciality and contractivity of $x^{1/(2r)}$.
The compressed $S'_n$ converges in norm to $S'_n$ as $r\to\infty$.
Lower semicontinuity yields
\begin{equation}\label{trace:localized-norm-error}
 \lambda(S'_n)\leq\lambda(S_n)+\Delta_{n,j}d_\tau(b).
\end{equation}
All of $\lambda(w_n)$, $\lambda(S_n)$, $\lambda(S'_n)$, and
$\tau(x)$ are finite in the case under consideration.  We may
therefore subtract these trace values.  Since
$x=S'_n+u_{n,0}$ and $w_n=S_n+u_{n,1}$,
\eqref{trace:retained-finite} and
\eqref{trace:localized-norm-error} give
\[
 \lambda(u_{n,1})
 \leq\lambda(u_{n,0})
       +(\delta_{n,j}+\Delta_{n,j})d_\tau(b).
\]
Together with \eqref{trace:first-small-part}, this proves, for both
$i=0,1$,
\[
 \lambda(u_{n,i})
 \leq(\delta_{n,j}+\Delta_{n,j}+\ell_n)d_\tau(b).
\]
Set
\begin{equation}\label{trace:error-coefficient}
 \varepsilon_{n,j}=\delta_{n,j}+\Delta_{n,j}+\ell_n.
\end{equation}
The trace-cut estimate \eqref{prelim:trace-cut} proves
\eqref{trace:small-rank}.  In particular, when $d_\tau(b)=0$ the
right side is exactly zero, not an absolute error term.

We now prove \eqref{trace:large-rank}.  Suppose $\ell_n\leq1/4$ and
write
\[
 t_n=(z_n-\gamma/16)_+.
\]
Because $(q_n-1/2)_+\geq1/2-e_n$ and
$t_n\geq z_n-(\gamma/16)1$, we have
\begin{equation}\label{trace:expectation-lower}
 L(t_n)\geq(1/2-\ell_n)a-(\gamma/16)1.
\end{equation}
On the support of $b$, the element $a$ is at least $\gamma/2$.
Consequently compression by $b^{1/(2r)}$ gives
\begin{equation}\label{trace:compressed-lower}
 b^{1/(2r)}L(t_n)b^{1/(2r)}
 \geq(\gamma/16)b^{1/r}.
\end{equation}
If $d_\tau(b)<\infty$, the compression
$b^{1/(2r)}t_nb^{1/(2r)}$ belongs to
$\her_A(b)\otimes\Z_1$, so
\eqref{trace:localized-expectation} and traciality imply
\[
 \lambda(t_n)
 \geq\tau(b^{1/(2r)}L(t_n)b^{1/(2r)})
 \geq(\gamma/16)\tau(b^{1/r}).
\]
Let $r\to\infty$ and use
$d_\lambda(t_n)\geq\lambda(t_n)$, since $t_n$ is a contraction.
This proves \eqref{trace:large-rank} in the finite-rank case,
including rank zero.

For the infinite-rank case we use ideals, without appealing to
\eqref{trace:localized-expectation}.  First,
\begin{equation}\label{trace:expectation-ideal}
 L(t_n)\otimes1\in\Ideal_R(t_n).
\end{equation}
To see this, apply the domination theorem,
Theorem~\ref{prelim:domination}, on $\Z_1$ to
$v\mapsto\tau_{\Z}(v)1$ and the identity map.  The source is exact,
the maps are nuclear, the identity has order zero, and the required
pointwise ideal containment follows from simplicity of $\Z$.
There are therefore completely positive maps of the form
\[
 V_l(v)=\sum_s z_{l,s}^*vz_{l,s}
\]
converging pointwise in norm to $\tau_{\Z}(v)1$.  Including $1$ in
the approximation tests allows us to require $\|V_l\|\leq2$.
Tensoring with $\id_A$ then preserves pointwise norm convergence:
it holds first on elementary tensors, and the uniform norm bound
extends it to their completion.  Each resulting approximation to
$L(t_n)\otimes1$ is a sum of compressions of $t_n$ in $R$, proving
\eqref{trace:expectation-ideal}.

Let $I=\Ideal_R(t_n)$.  In $R/I$, equation
\eqref{trace:expectation-lower} and
\eqref{trace:expectation-ideal} give
\[
 0\geq(1/2-\ell_n)(a+I)-(\gamma/16)(1+I).
\]
Thus $\|a+I\|\leq\gamma/4$, and $(a-\gamma/4)_+\in I$.
Taking one more positive cut gives
\[
 b=\bigl((a-\gamma/4)_+-\gamma/4\bigr)_+,
 \qquad [b]\leq M[t_n]
\]
for some finite positive integer $M$.  If $d_\tau(b)=\infty$, then
$d_\lambda(b)=\infty$, so this comparison forces
$d_\lambda(t_n)=\infty$.  This finishes the proof in all cases.
\end{proof}

\subsection{Slowly chosen corners and bounded comparison witnesses}

We let the number of available orthogonal copies in the second
Jiang--Su factor grow slowly with the coordinate, so that the
localized comparisons remain valid.

\begin{proposition}[Localized comparison]\label{trace:comparison}
There is a sequence $m(n)\geq2$ with $m(n)\to_\omega\infty$ such
that, on putting
\[
 P_n=D_{m(n)}\in\Z_2,\qquad P=(P_n)_\omega\in B,
 \qquad
 \Psi(a)=P(q-1/2)_+\pi(a),
\]
the map $\Psi\colon A\to B$ is c.p.c. order zero and, for every
$j\geq1$,
\begin{equation}\label{trace:small-part-comparison}
 e\pi(x_j),\ e\theta(x_j)
 \precsim (\Psi(a_j)-c_j)_+
 \quad\text{in }B,
 \qquad c_j=\gamma_j/64.
\end{equation}
\end{proposition}

\begin{proof}
Temporarily fix $m\geq2$ and $j$, and keep the abbreviations from
the preceding proof.  The copy of $D_m$ is now in $\Z_2$, separate
from $z_n$ and $t_n=(z_n-\gamma/16)_+$.  The compression identity
in \eqref{supplies:copy-relations} gives
\begin{equation}\label{trace:separate-leg-comparison}
 [t_n]\leq K_m[t_n(D_m-1/2)_+]
       \leq K_m[(z_nD_m-\gamma/32)_+]
 \quad\text{in }W(R).
\end{equation}
For the second inequality, use joint functional calculus and the
scalar inequality
\[
 (s-\gamma/16)_+(t-1/2)_+
      \leq(st-\gamma/32)_+\qquad(0\leq s,t\leq1).
\]
Where the left side is positive, the difference between the
uncut expression on the right and the left side is
$s/2+(\gamma/16)t-\gamma/16\geq0$; elsewhere the assertion is
immediate.

Enumerate the countable set of pairs $(j,\eta)$ with
$j\geq1$ and $\eta\in\mathbb Q_{>0}$.  For each $m\geq2$ let
$\Omega_m$ be an $\omega$-large set on which $\ell_n\leq1/4$ and
\begin{equation}\label{trace:slow-requirement}
 \eta^{-1}\varepsilon_{n,j}
       \leq\frac{\gamma_j}{32K_m}
\end{equation}
for all the first $m$ pairs.  Such a set exists because the finitely
many error coefficients tend to zero along $\omega$.  Define
\[
 E_m=\{n:n\geq m\}\cap\bigcap_{r=2}^m\Omega_r.
\]
The sets $E_m$ are decreasing and $\omega$-large.  If $n$ belongs
to at least one $E_m$, let $m(n)$ be the largest such $m$; otherwise
put $m(n)=2$.  The largest value exists since membership in $E_m$
requires $m\leq n$.  This construction gives
$m(n)\to_\omega\infty$ and ensures that, for every fixed pair
$(j,\eta)$, condition \eqref{trace:slow-requirement} with
$m=m(n)$ holds on an $\omega$-large set.

Set $P_n=D_{m(n)}$.  For each fixed $(j,\eta)$, on the preceding
$\omega$-large set and for every lower semicontinuous extended
trace $\lambda$ on $R$, we claim
\begin{equation}\label{trace:uniform-gap}
 d_\lambda((u_{n,i}-\eta)_+)
 \leq\frac12 d_\lambda((z_nP_n-\gamma_j/32)_+)
 \qquad(i=0,1).
\end{equation}
When $d_\tau(b_j)<\infty$, combine
\eqref{trace:small-rank}, \eqref{trace:large-rank},
\eqref{trace:separate-leg-comparison}, and
\eqref{trace:slow-requirement}:
\[
 \begin{split}
 d_\lambda((u_{n,i}-\eta)_+)
 &\leq \frac{\gamma_j}{32K_{m(n)}}d_\tau(b_j)\\
 &\leq\frac12
      d_\lambda((z_nP_n-\gamma_j/32)_+).
 \end{split}
\]
This includes $d_\tau(b_j)=0$.  If $d_\tau(b_j)=\infty$, then
\eqref{trace:large-rank} gives $d_\lambda(t_n)=\infty$, and
\eqref{trace:separate-leg-comparison}, with finite $K_{m(n)}$,
forces the right side of \eqref{trace:uniform-gap} to be infinite.
Thus \eqref{trace:uniform-gap} also holds in that case.  Its
coordinate set is independent of $\lambda$.

The element $P$ commutes with the entire first subsystem $B_1$.
Also $q$ commutes with $\pi(A)$.  Hence $\Psi$ is a homomorphism
$\pi$ multiplied by commuting positive contractions, and is c.p.c.
order zero.  In $B$ we have
\begin{equation}\label{trace:coordinate-identifications}
 (u_{n,0})_\omega=e\pi(x_j),\qquad
 (u_{n,1})_\omega=e\theta(x_j),\qquad
 (z_nP_n)_\omega=\Psi(a_j).
\end{equation}

Since $R$ is exact and $\Z$-stable,
Lemma~\ref{prelim:half-comparison} applies to
\eqref{trace:uniform-gap}.  Thus, on an $\omega$-large set,
\[
 (u_{n,i}-\eta)_+
   \precsim(z_nP_n-2c_j)_+.
\]
We explain the passage to the ultrapower, because the norm bound
on witnesses is essential.  Fix $j$, a positive rational $\eta$,
$i\in\{0,1\}$, and an accuracy $\epsilon>0$, and let
$u_i=(u_{n,i})_\omega$.  By
Lemma~\ref{prelim:bounded-witness}, choose witnesses $r_n\in R$
on this large set such that
\[
 \begin{split}
 \bigl\|r_n^*(z_nP_n-c_j)_+r_n-(u_{n,i}-\eta)_+\bigr\|
     &<\epsilon,\\
 \|r_n\|^2&\leq(1+\epsilon)/c_j.
 \end{split}
\]
Set $r_n=0$ elsewhere.  This is a bounded sequence, so it defines
$r\in B$.  By \eqref{trace:coordinate-identifications},
\[
 \bigl\|r^*(\Psi(a_j)-c_j)_+r
       -(u_i-\eta)_+\bigr\|\leq\epsilon.
\]
As $\epsilon$ is arbitrary, each displayed source cut is Cuntz
below $(\Psi(a_j)-c_j)_+$.  Letting positive rational $\eta$
decrease to zero proves \eqref{trace:small-part-comparison}.
\end{proof}

\begin{lemma}\label{trace:copies}
For every positive integer $L$ there are contractions
$v_1,\ldots,v_L\in B$, represented entirely in the $\Z_2$ leg,
such that
\begin{equation}\label{trace:ultrapower-copies}
 v_l^*v_{l'}=\delta_{ll'}P\qquad(1\leq l,l'\leq L).
\end{equation}
In particular these elements commute with $B_1$, and hence with
$\pi(A)$ and $q$.
\end{lemma}

\begin{proof}
At coordinates where $m(n)\geq L$, use the elements
$v_{m(n),l}$ from Lemma~\ref{supplies:copies}; at all other
coordinates use zero.  The representatives have norm at most one.
Since $m(n)\to_\omega\infty$, their relations in
\eqref{supplies:copy-relations} give
\eqref{trace:ultrapower-copies}.  Their tensor-factor location gives
the stated commutation.
\end{proof}

\section{The supporting map and its relative commutant}
\label{sec:commutant}

We retain the notation of the preceding section.  In particular,
\[
 R=A\otimes\Z_1\otimes\Z_2\otimes\Z_3,
 \qquad B=R_\omega,
\]
where $A$ is nonzero, separable, unital and nuclear.  The maps
$\pi,\theta\colon A\to B$ have ranges in the first subsystem
$(A\otimes\Z_1)_\omega$, with $\pi$ the constant first-factor
homomorphism.  The positive contractions $e,q=1-e$ satisfy
\eqref{setup:common-equations}.  The second-leg contraction $P$ defines
\[
 \Psi(a)=P(q-1/2)_+\pi(a),\qquad a\in A.
\]
We shall use both the cut comparisons \eqref{trace:small-part-comparison} and
the arbitrarily long finite families of second-leg elements in
\eqref{trace:ultrapower-copies}.  The constant third leg will remain available
for the functional-calculus argument.

\subsection{A supporting map preserving the relevant ideals}

Put $h=\theta(1_A)$.  An abstract supporting homomorphism for an
order-zero map need not take values in its codomain.  We therefore seek
a c.p.c. order-zero map $\widehat\theta\colon A\to B$ with three
roles.  Its factorization $h\widehat\theta=\theta$ recovers the model
and gives $h\widehat\theta(1_A)=h$, the support identity for $h$ in
the second matrix corner below.  The equality
$q\widehat\theta=q\pi$ preserves the common supported part.
Finally, pointwise ideal containment of $e\widehat\theta(a')$ and
$e\pi(a')$ in $\Ideal_B(\Psi(a'))$ supplies the hypothesis for
Gabe's theorem, Theorem~\ref{prelim:domination}, to approximate the
error maps by finite sums of compressions of $\Psi$.

\begin{proposition}[A supporting map preserving ideals]\label{support:map}
There is a c.p.c. order-zero map $\widehat\theta\colon A\to B$
whose range commutes with the constant copy of $\Z_3$ and such that
\begin{align}
 h\widehat\theta(a)=\widehat\theta(a)h&=\theta(a),
       \qquad a\in A,\label{support:factorization}\\
 [e,\widehat\theta(a)]&=0,
 \qquad q\widehat\theta(a)=q\pi(a),\qquad a\in A,
       \label{support:common}\\
 e\widehat\theta(a'),\ e\pi(a')&\in\Ideal_B(\Psi(a')),
       \qquad a'\in A_+.
       \label{support:ideals}
\end{align}
\end{proposition}

\begin{proof}
Consider the order-zero root maps $\theta^{1/l}$, for $l\geq1$.
They commute with $e$ and the constant third leg.  Indeed, each value
$\theta^{1/l}(a)$ is a norm limit of sandwiches of $\theta(a)$ by
bounded continuous functions of $h$.  One obtains these sandwiches
by truncating $t^{(1/l-1)/2}$ from above; the resulting scalar weights
converge uniformly to $t^{1/l}$ after multiplication by $t$.
The same observation gives
\begin{equation}\label{support:root-comparison}
 e\theta^{1/l}(x_j)\precsim e\theta(x_j).
\end{equation}
Here $e$ commutes with the values involved, so the sandwiches also
approximate the positive element on the left.

Taking $a=1_A$ in \eqref{setup:common-equations} gives $qh=q$.  Continuous
functional calculus therefore gives $qh^{1/l}=q$, and the sandwich
formula together with \eqref{setup:common-equations} yields
\[
 q\theta^{1/l}(a)=q\pi(a),\qquad a\in A.
\]
Moreover,
\[
 h\theta^{1/l}(a)=\theta^{1/l}(a)h\longrightarrow\theta(a)
 \quad(l\longrightarrow\infty),
\]
by the uniform convergence of $t^{1+1/l}$ to $t$ on $[0,1]$.

For every fixed $j$, combining \eqref{support:root-comparison} with
\eqref{trace:small-part-comparison} gives
\[
 e\theta^{1/l}(x_j)\precsim(\Psi(a_j)-c_j)_+,
 \qquad c_j=\gamma_j/64>0.
\]
The bounded-witness estimate of Lemma~\ref{prelim:bounded-witness}, with the
spare cut $c_j$, now shows that this positive element can be
approximated arbitrarily closely by $r_j^*\Psi(a_j)r_j$, with
\begin{equation}\label{support:witness-bound}
 \|r_j\|^2\leq 2/c_j.
\end{equation}
The bound depends on $j$ but, crucially, not on $l$.

We apply the countable diagonalization principle of
Lemma~\ref{prelim:epsilon-test}.  To make the boundedness issue explicit,
lift each root map to a sequence of c.p.c. maps into $R$ by the
Choi--Effros lifting theorem.  Keep fixed representatives of
$e,q,h,\theta,\Psi$ and $\pi$.  The tests impose the factorization
and common-part identities on countable dense subsets of $A$,
commutation with $e$ and a dense subset of the constant third leg,
and order zero on a countable dense family of orthogonal pairs of
positive contractions.  For each $j$ include a witness variable with
a fixed representative bound strictly larger than
$\sqrt{2/c_j}$, and impose
\[
 e\widehat\theta(x_j)=r_j^*\Psi(a_j)r_j.
\]
Every finite list is approximately solvable: take $l$ sufficiently
large for the factorization tests, and then use the witnesses above
with sufficiently small errors.  Their fixed bounds allow the
ultrapower tests to be realized by bounded coordinate variables.
The diagonalization produces a c.p.c. map satisfying all the tests.
Continuity extends the identities and commutations to all the
specified elements; the dense orthogonality tests make the map
order zero.  This proves \eqref{support:factorization} and
\eqref{support:common}, as well as the witness identities for all $j$.

It remains to extend ideal containment from the listed cuts in
\eqref{model:tests} to an arbitrary positive contraction $a'\in A$.
Such an $a'$ can be
approximated arbitrarily closely by $x_j$ with
$a_j\in\Ideal_A(a')$.  To see this, choose $d$ in the fixed dense
positive family close to $a'$, and a small positive rational
$r>\|d-a'\|$.  The image of $d$ in $A/\Ideal_A(a')$ has norm less
than $r$, so $(d-r)_+\in\Ideal_A(a')$.  This cut belongs to the
chosen family; take it as $a_j$ and choose the accompanying rational
$\gamma_j$ small.  Then $x_j=(a_j-\gamma_j)_+$ is as close to $a'$
as desired.

The map $\Psi$ is a positive weight times the homomorphism $\pi$,
and the weight commutes with $\pi(A)$.  Consequently
\[
 a_j\in\Ideal_A(a')
 \quad\Longrightarrow\quad
 \Psi(a_j)\in\Ideal_B(\Psi(a')):
\]
apply $\Psi$ to finite sums of products through $a'$ and pass to
norm limits.  The witness identity places
$e\widehat\theta(x_j)$ in this ideal.  The comparison
\eqref{trace:small-part-comparison} similarly places $e\pi(x_j)$ there.
Taking norm limits, and then scaling, proves
\eqref{support:ideals} for every $a'\in A_+$.
\end{proof}

\subsection{Comparison inside the supported relative commutant}

Define the c.p.c. order-zero map
\[
 \rho=\operatorname{diag}(\rho_0,\rho_1)\colon A\longrightarrow M_2(B),
 \qquad \rho_0=\pi,\quad\rho_1=\widehat\theta,
\]
and its supported relative commutant
\begin{equation}\label{comm:definition}
 C=M_2(B)\cap\rho(A)'\cap\{1-\rho(1_A)\}^{\perp}.
\end{equation}
The last condition means annihilation on both sides.  A matrix in $C$
supported in position $(0,1)$ has an entry that intertwines $\rho_1$
with $\rho_0$ and satisfies both support conditions.  The matrix's two
modulus squares lie in the diagonal corners of $C$ and have the same Cuntz
class, allowing comparison of positive elements in the two corners
inside one semigroup.  For $b\in B$,
write $b^{(i)}$ for the matrix with entry $b$ in position $(i,i)$ and
zeros elsewhere, using indices $i=0,1$.  Put
\[
 C_i=\{b\in B:b^{(i)}\in C\},\qquad
 Q_i=q^{(i)},\qquad H=1_B^{(0)}.
\]
The common-part identities show that $Q_i,H\in C$.  They also imply
that the matrix with $q^{1/2}$ in position $(0,1)$ belongs to $C$:
positive-support functional calculus passes the intertwining and
support identities from $q$ to $q^{1/2}$.  Its two modulus squares
are $Q_1$ and $Q_0$.  Thus
\begin{equation}\label{comm:common-equivalence}
 [Q_0]=[Q_1]\quad\text{in }W(C).
\end{equation}

\begin{proposition}[Packing the error part]\label{comm:packing}
For $i\in\{0,1\}$ and every positive contraction $c\in C_i$,
\begin{align}
 m[(c^{1/2}ec^{1/2})^{(i)}]&\leq[Q_0],
       &&m\in\N,\label{comm:small-part}\\
 [c^{(i)}]&\leq2[Q_0]
       &&\text{in }W(C).\label{comm:two-common}
\end{align}
\end{proposition}

\begin{proof}
Fix $m\geq1$ and set $t=c^{1/2}ec^{1/2}$.  This positive
contraction commutes with $\rho_i(A)$ and satisfies
$t\rho_i(1_A)=t$.  The order-zero product identity therefore gives
\begin{equation}\label{comm:supported-product}
 t\rho_i(a)\rho_i(b)=t\rho_i(ab),\qquad a,b\in A.
\end{equation}
Apply Theorem~\ref{prelim:domination} to
\[
 \Phi(a)=t\rho_i(a),\qquad
 \Psi(a)=P(q-1/2)_+\pi(a).
\]
Both maps are completely positive and nuclear, since their domain
is nuclear, and $\Psi$ is order zero.  For $a\in A_+$,
\[
 \Phi(a)=c^{1/2}e\rho_i(a)c^{1/2}
       \in\Ideal_B(\Psi(a))
\]
by Proposition~\ref{support:map}.  Hence finite sums
$\sum_{l=1}^L y_l^*\Psi(\,\cdot\,)y_l$ approximate $\Phi$
pointwise in norm.

Write $h_q=(q-1/2)_+$.  For one such approximation, take $mL$
second-leg elements $v_{sl}$ from \eqref{trace:ultrapower-copies} and set
\[
 r_s=\sum_{l=1}^L v_{sl}h_q^{1/2}y_l,
 \qquad 1\leq s\leq m.
\]
These second-leg elements commute with $q$ and $\pi(A)$, and satisfy
$v_{sl}^*v_{s'l'}=\delta_{ss'}\delta_{ll'}P$.  Expanding the products
therefore gives
\begin{equation}\label{comm:packed-approximation}
 r_s^*\pi(a)r_{s'}=
 \delta_{ss'}\sum_{l=1}^L y_l^*\Psi(a)y_l,
 \qquad a\in A.
\end{equation}
For $f(u)=\min\{1,2u\}$ on $[0,1]$ one also has
$f(q)r_s=r_s$, since $f(q)h_q^{1/2}=h_q^{1/2}$.
Include $1_A$ in the approximation set.  As
$\Phi(1_A)=t$ and $\|t\|\leq1$, sufficiently accurate
approximations in \eqref{comm:packed-approximation} have
$\|r_s\|^2\leq2$.  Thus, for this fixed $m$, there is a fixed
bound on every packed variable, independent of $L$ and of the
approximation accuracy.

Using bounded coordinate lifts and the diagonalization principle of
Lemma~\ref{prelim:epsilon-test} on a countable dense subset of $A$, we
obtain $r_1,\ldots,r_m\in B$ with
\begin{equation}\label{comm:packed-exact}
 f(q)r_s=r_s,\qquad
 r_s^*\pi(a)r_{s'}=\delta_{ss'}t\rho_i(a),
 \qquad a\in A.
\end{equation}
These elements are intertwiners.  Indeed, using
\eqref{comm:supported-product} and
\eqref{comm:packed-exact}, the right modulus square of
$\pi(a)r_s-r_s\rho_i(a)$ is
\[
 t\rho_i(a^*a)-t\rho_i(a^*)\rho_i(a)=0.
\]
Consequently $\pi(a)r_s=r_s\rho_i(a)$ for all $a\in A$.
Taking $a=1_A$ gives $r_s\rho_i(1_A)=r_s$ as well.
The matrix $R_s$ with $r_s$ in position $(0,i)$ thus belongs to
$C$, including its support condition.

For the row $V=(R_1,\ldots,R_m)$ over $C$, the positive matrix
$V^*V$ has $m$ diagonal entries equal to $t^{(i)}$ and all
other entries zero.  Meanwhile $VV^*\in\her_C(Q_0)$:
\eqref{comm:packed-exact} gives
$f(Q_0)VV^*f(Q_0)=VV^*$, and $f(0)=0$.
Equality of the Cuntz classes of the two modulus squares now proves
\eqref{comm:small-part}.

Finally,
\[
 c=c^{1/2}ec^{1/2}+c^{1/2}qc^{1/2}.
\]
The second summand, placed in corner $i$, is Cuntz below $Q_i$
inside $C$.  Apply \eqref{comm:small-part} with $m=1$, the
subadditivity of Cuntz classes of positive sums, and
\eqref{comm:common-equivalence}.  This proves
\eqref{comm:two-common}.
\end{proof}

\begin{corollary}\label{comm:fullness}
The elements $Q_0,Q_1$ and $H$ are full in $C$, and $C$ is full
in $M_2(B)$.
\end{corollary}

\begin{proof}
Every positive element of $M_2(B)$ is bounded by twice its diagonal
part.  For an element of $C_+$, each diagonal corner also lies in
$C$, because $\rho$ and its support condition are diagonal.
After scaling, \eqref{comm:two-common} shows that these corners
belong to the ideal generated by $Q_0$.  Hence so does the original
positive element.  Thus $Q_0$ is full; the same holds for $Q_1$ by
\eqref{comm:common-equivalence}, and for $H$ because $Q_0\leq H$.
The corner projection $H$ is full in the ambient matrix algebra
$M_2(B)$, so the latter also contains $C$ as a full subalgebra.
\end{proof}

\begin{proposition}\label{comm:unperforated}
The ordered semigroup $W(C)$ is almost unperforated.
\end{proposition}

\begin{proof}
Suppose $(N+1)x\leq Ny$ in $W(C)$, with $N\geq1$.
Choose a separable $C^*$-subalgebra $C_*\subset C$ containing the
matrix entries of representatives of $x,y$ and of countably many
comparison witnesses.  The inequality then already holds over
$C_*$.

There is a unital copy of $\Z$ in $M_2(B)$ commuting with both
$C_*$ and $\rho(A)$.  Here is the coordinate argument, which also
allows the representatives of $C_*$ to vary with the coordinate.
Choose isomorphisms
$\sigma_l\colon\Z\otimes\Z\to\Z$ with
$\sigma_l(d\otimes1)\to d$, and put
$\nu_l(z)=\sigma_l(1\otimes z)$.  For fixed $d\in\Z$,
\[
 \|[\nu_l(z),d]\|
 \leq 2\|z\|\,\|d-\sigma_l(d\otimes1)\|.
\]
Thus the commutators tend to zero uniformly for $z$ in the unit
ball.  Place these embeddings diagonally on the third tensor leg
of $M_2(R)$.  Tensor approximation gives the same uniform
commutation property with any fixed element of $M_2(R)$.
At coordinate $n$, choose $l$ sufficiently large for the entries
of representatives of the first $n$ elements of countable dense
commutation data, with tolerance tending to zero.  The resulting
coordinate embeddings induce a unital embedding
$\nu\colon\Z\to M_2(B)$ commuting with $C_*$ and $\rho(A)$.
It is isometric because each coordinate embedding is isometric.

Multiplication gives a homomorphism
\[
 C_*\otimes\Z\longrightarrow C,\qquad c\otimes z\longmapsto c\nu(z).
\]
Commutation gives a homomorphism on the maximal tensor product;
nuclearity of $\Z$ identifies it with the spatial tensor product.
The range satisfies the support condition because $c$ does and
$\nu(z)$ commutes with $\rho(1_A)$.  This homomorphism restricts
to the original inclusion on $C_*\otimes1$.

The algebra $C_*\otimes\Z$ is $\Z$-stable.  Its Cuntz semigroup
is therefore almost unperforated by Theorem~\ref{prelim:comparison}.  Tensoring the original inequality with
$1_\Z$ gives comparison of the images of $x$ and $y$ there.
Mapping the comparison back to $C$ yields $x\leq y$.

Here the reindexed copy $\nu(\Z)$ is used for this inequality alone.
The constant third-leg copy retained above is the separate source of
the weight $k$ below; no commutation between these two copies is
required.
\end{proof}

We next use the full projection $H$, the corner divisibility supplied
by Lemma~\ref{supplies:copies}, and almost unperforation of $W(C)$
to verify the structural hereditary hypothesis of the BBSTWW
criterion recalled in Lemma~\ref{prelim:unitary}.

\begin{proposition}\label{conjugacy:hereditary}
Let $D\subset C$ be a full hereditary subalgebra.  Suppose that
$y\in D$ and that totally full positive elements $e_l,e_r\in D$
satisfy
\[
 e_l y=ye_r=0.
\]
Then there is a full element $s\in D$ with $sy=ys=0$.
\end{proposition}

\begin{proof}
Fullness of $e_l,e_r$ in $D$ and fullness of $D$ in $C$ imply
that both elements are full in $C$.  Since $H$ is a projection,
there is a common finite integer $N\geq1$ such that
\begin{equation}\label{conjugacy:projection-domination}
 [H]\leq N[e_l],\qquad [H]\leq N[e_r]
 \quad\text{in }W(C).
\end{equation}
Use the positive contraction $D_{N+1}$ from the order-zero supplies
on a constant $\Z$ leg, and put
$p_0=(D_{N+1})^{(0)}$.  It belongs to $C$, since the first-corner
map is $\pi$, and it has norm one.  Simplicity of the constant
leg gives fullness of $p_0$ in the corner with unit $H$; because
$H$ is full in $C$, so is $p_0$.

The $N+1$ accompanying elements from Lemma~\ref{supplies:copies},
placed in this corner, form a row over $C$ whose right modulus
square is the diagonal sum of $N+1$ copies of $p_0$, while its
left modulus square belongs to $HCH$.  Hence
\[
 (N+1)[p_0]\leq[H].
\]
Combining this with \eqref{conjugacy:projection-domination} and
Proposition~\ref{comm:unperforated} gives
$p_0\precsim e_l,e_r$ in $C$.  We may therefore choose
$w_l,w_r\in C$ so that
\[
 A_l=w_l^*e_lw_l,\qquad A_r=w_r^*e_rw_r
\]
are as close to $p_0$ as prescribed.

We specify the accuracy before choosing these witnesses.
Fullness of $p_0^2$ in the constant simple leg gives fixed
$z_1,\ldots,z_L\in HCH$ with
\[
 \sum_{j=1}^L z_j^*p_0^2z_j=H.
\]
Choose $\varepsilon>0$ such that
\begin{equation}\label{conjugacy:fullness-tolerance}
 (2\varepsilon+\varepsilon^2)
       \sum_{j=1}^L\|z_j\|^2<1,
\end{equation}
and choose $w_l,w_r$ with $\|A_l-p_0\|,\|A_r-p_0\|<\varepsilon$.
Set
\[
 s=e_r^{1/2}w_rw_l^*e_l^{1/2}.
\]
Heredity gives $s\in D$, since $DCD\subset D$.
Functional calculus applied to the annihilation identities gives
$e_l^{1/2}y=ye_r^{1/2}=0$, so $sy=ys=0$.

Let $J=\Ideal_C(s)$.  It contains
\[
 w_r^*e_r^{1/2}s\,e_l^{1/2}w_l=A_rA_l,
 \qquad
 \|A_rA_l-p_0^2\|\leq2\varepsilon+\varepsilon^2.
\]
Thus the element $b=\sum_j z_j^*A_rA_lz_j\in J$ satisfies
$\|H-b\|<1$ by \eqref{conjugacy:fullness-tolerance}.
The image of $H$ in $C/J$ is a projection of norm less than one,
and is therefore zero.  Since $H$ is full in $C$, this forces
$J=C$.

It remains to prove that $s$ is full inside $D$.  Let
$I=\Ideal_D(s)$.  For any ideal $I$ of a hereditary
subalgebra $D\subset C$, one has
\[
 D\,\Ideal_C(I)\,D\subset I.
\]
Indeed, the ideal in the middle is the closed linear span of
products through elements of $I$.  By linear decomposition it is
enough to use a positive middle element $i\in I_+$, and then
\[
 dci c'd'
   =(dc\,i^{1/3})\,i^{1/3}\,(i^{1/3}c'd')\in I:
\]
the two outer factors lie in $D$ by heredity.
Here $\Ideal_C(I)$ contains $J=C$, so the displayed inclusion
and an approximate identity of $D$ give $D\subset I$.
Hence $I=D$, proving that $s$ is full in $D$.
\end{proof}

\subsection{Full-spectrum weights and unitary equivalence}

The structural conditions on $C$ in Lemma~\ref{prelim:unitary} are now
verified.  We turn to the remaining input: totally full contractions
with matching Cuntz classes under positive functional calculus.

Fix a positive contraction $k$ of spectrum $[0,1]$ in the retained
constant third leg.  Such a $k$ is obtained from the scalar interval
coordinate in a prime dimension-drop algebra embedded in $\Z_3$.
When necessary, a third-leg element also denotes its diagonal
copy in $M_2(B)$.  Set
\begin{equation}\label{conjugacy:weights}
 h_0=k^{(0)},\qquad h_1=(kh)^{(1)}.
\end{equation}
These are positive contractions in $C$.  In the second corner,
the support condition follows from
$h\widehat\theta(1_A)=h$; commutation follows from
Proposition~\ref{support:map}.

\begin{proposition}\label{conjugacy:functions}
The elements $h_0,h_1$ are totally full in $C$.  For every
$g\in C_0((0,1])_+$, the elements $g(h_0)$ and $g(h_1)$ are
Cuntz equivalent in $C$.
\end{proposition}

\begin{proof}
The case $g=0$ is immediate.  Otherwise rescale so that
$\|g\|\leq1$, and put
\[
 G_i=g(h_i),\qquad T_i=Q_iG_i,\qquad E_i=G_i-T_i=e^{(i)}G_i.
\]
All these elements are positive and belong to $C$.  The factors
commute because $e$ commutes with $h$ and is represented in the
first subsystem.  Since $qh=q$, functional calculus gives
$T_i=Q_i g(k)$.  The matrix with
$q^{1/2}g(k)^{1/2}$ in position $(0,1)$ belongs to $C$ and has
modulus squares $T_1,T_0$.  Consequently
\begin{equation}\label{conjugacy:common-functions}
 [T_0]=[T_1]\quad\text{in }W(C).
\end{equation}

The positive element $g(k)$ is nonzero in the simple unital
constant third leg.  Choose finitely many elements
$z_1,\ldots,z_K$ in that leg with
\[
 \sum_{l=1}^K z_l^*g(k)z_l=1.
\]
For $\delta>0$ let
\[
 b_\delta(u)=
 \begin{cases}
 ((u-\delta)_+/u)^{1/2},&u>0,\\
 0,&u=0.
 \end{cases}
\]
The diagonalized $z_l$ commute with $Q_i$ and $\rho(A)$ and
multiply $C$ into itself.  In particular
$z_l b_\delta(Q_i)\in C$, and
\[
 (Q_i-\delta)_+
 =\sum_{l=1}^K
   (z_l b_\delta(Q_i))^*T_i(z_l b_\delta(Q_i)).
\]
Taking all cuts proves
\begin{equation}\label{conjugacy:fullness-bound}
 [Q_i]\leq K[T_i]\leq K[G_i]\quad\text{in }W(C).
\end{equation}
Since $Q_i$ is full, this shows that $g(h_i)$ is full for every
nonzero $g$.  In particular the two elements in
\eqref{conjugacy:weights} are totally full; the same argument
also ensures that every such functional calculus is nonzero.

Let $d$ be any dimension function on $W(C)$ normalized at either
$[G_0]$ or $[G_1]$, if one exists.  Equations
\eqref{comm:two-common} and
\eqref{conjugacy:fullness-bound}, together with
\eqref{comm:common-equivalence}, imply
\[
 0<d([Q_0])=d([Q_1])<\infty.
\]
Apply \eqref{comm:small-part} with $c^{(i)}=G_i$.
Since $G_i$ commutes with $e^{(i)}$, it follows that
$m[E_i]\leq[Q_0]$ for every $m$, and hence $d([E_i])=0$.
The inequalities
\[
 [T_i]\leq[G_i]\leq[T_i]+[E_i]
\]
and \eqref{conjugacy:common-functions} now give
\begin{equation}\label{conjugacy:dimension-equality}
 d([G_0])=d([G_1])=1.
\end{equation}

Fix $\eta>0$.  By continuity, $g(0)=0$, and $g\ne0$, there is
a nonzero continuous function $g'$ with $0\leq g'\leq g$ and
support contained in the open set where $0<g<\eta$.
For the dimension functions under consideration,
\eqref{conjugacy:fullness-bound} applied to $g'$ gives
$d([g'(h_i)])>0$.  This element is orthogonal to
$(G_i-\eta)_+$, and their sum is bounded by $G_i$.
Therefore \eqref{conjugacy:dimension-equality} gives
\[
 d([(G_i-\eta)_+])+d([g'(h_i)])\leq1,
 \qquad
 d([(G_i-\eta)_+])<1=d([G_{1-i}]).
\]
This holds for every dimension function normalized at
$[G_{1-i}]$.  The target $G_{1-i}$ is full, so the ideal condition
for the comparison criterion in Theorem~\ref{prelim:comparison} holds.
Proposition~\ref{comm:unperforated} and that criterion yield
$(G_i-\eta)_+\precsim G_{1-i}$.  This use of the criterion
also permits an empty family of normalized dimension functions.
Taking all positive cuts, and then exchanging $i=0,1$, proves
the required Cuntz equivalence.
\end{proof}

\begin{proposition}[Weighted unitary equivalence]\label{conjugacy:weighted}
Put $k^{[0]}=k$ and $k^{[1]}=1-k$, in the constant third leg.
There are unitaries $U_0,U_1\in M_2(B)$ such that
\begin{equation}\label{conjugacy:weighted-equation}
 U_i\bigl(k^{[i]}\theta(a)\bigr)^{(1)}U_i^*
   =\bigl(k^{[i]}\pi(a)\bigr)^{(0)},
 \qquad a\in A,\quad i=0,1.
\end{equation}
They can be chosen to commute with $\rho(A)$.
\end{proposition}

\begin{proof}
Apply the relative-commutant unitary-equivalence criterion of
Lemma~\ref{prelim:unitary}, namely \cite[Lemma~5.3]{BBSTWW}, to
\eqref{comm:definition}.  Its coordinate algebras are $M_2(R)$,
which are separable, unital and $\Z$-stable, and its domain is the
separable unital algebra $A$.  The map $\rho$ is c.p.c. order
zero.  Corollary~\ref{comm:fullness} verifies fullness of $C$ in
the ambient ultrapower; Proposition~\ref{conjugacy:hereditary}
verifies the required condition in every full hereditary
subalgebra.  Proposition~\ref{conjugacy:functions} supplies total
fullness and all the functional-calculus comparisons.  Thus there
is a unitary in the unitization of $C$ conjugating $h_1$ to $h_0$.

Realizing this unitary in the ambient algebra as explained after
Lemma~\ref{prelim:unitary} gives $U_0\in M_2(B)$ commuting with
$\rho(A)$. Multiplying
$U_0h_1U_0^*=h_0$ by $\rho(a)$ and using
\eqref{support:factorization} gives
\[
 U_0(k\theta(a))^{(1)}U_0^*=(k\pi(a))^{(0)}.
\]
The element $1-k$ also has spectrum $[0,1]$ in the constant
third leg.  Repeating the same argument with this element, with
all the previously constructed maps unchanged, gives $U_1$ and
the other identity in \eqref{conjugacy:weighted-equation}.
\end{proof}

Complementary weights are part of the two-coloured approximation
framework developed in \cite{SWW,BBSTWW,CETWW}. Here the required
ideal-sensitive comparisons have been proved above. The two identities
in \eqref{conjugacy:weighted-equation} add to the first-corner inclusion
because $k+(1-k)=1$.  The next section
passes from these ultrapower identities to finite-dimensional
approximations with exactly two order-zero return maps.

\Needspace{22\baselineskip}
\section{Two colours and hereditary permanence}
\label{sec:endgame}

We first turn the ultrapower conjugacies into finite-dimensional
approximations of a corner inclusion. We then use stabilization to
retain the same two colours when the target is the algebra itself.

\begin{proposition}[Two-colour corner approximation]
\label{end:corner}
Let $A$ be a nonzero separable unital nuclear $C^*$-algebra and put
$R=A\otimes\Z^{\otimes3}$. For every finite $\mathcal F\subset A$
and $\varepsilon>0$, the map
\[
 \kappa_A:A\longrightarrow M_2(R),\qquad
 \kappa_A(a)=\begin{pmatrix}a\otimes1&0\\0&0\end{pmatrix}
\]
admits a finite-dimensional approximation with two colours:
there are a finite-dimensional algebra $F$, a c.p.c. map
$\alpha:A\to F$, and c.p.c. order-zero maps
$\varphi_0,\varphi_1:F\to M_2(R)$ such that
\[
 \|\varphi_0\alpha(a)+\varphi_1\alpha(a)-\kappa_A(a)\|
 <\varepsilon\qquad(a\in\mathcal F).
\]
\end{proposition}

\begin{proof}
Use the notation and the two unitaries of
Proposition~\ref{conjugacy:weighted}. Write $k^{[0]}=k$ and $k^{[1]}=1-k$.
Each unitary $U_i\in M_2(R_\omega)$ has unitary representatives
$U_{i,n}\in M_2(R)$. Indeed, a bounded representative is asymptotically
unitary on both sides. On an $\omega$-large set its coordinates are
invertible and polar correction gives unitaries with the same quotient
class; put the identity on the remaining coordinates.

For the matrix-model maps $\alpha_n:A\to F_n$ and
$\beta'_n:F_n\to A\otimes\Z_1\subset R$, define
\[
 \varphi_{i,n}(b)
 =U_{i,n}\begin{pmatrix}0&0\\0&k^{[i]}\beta'_n(b)\end{pmatrix}U_{i,n}^*.
\]
Each map is c.p.c. order zero: its weight is a positive contraction
in the third tensor factor, commuting with the range of $\beta'_n$;
corner inclusion and unitary conjugation preserve these properties.
The conjugacies in Proposition~\ref{conjugacy:weighted} give, for every $a\in A$,
\[
 \lim_{n\to\omega}
 \|\varphi_{0,n}\alpha_n(a)+\varphi_{1,n}\alpha_n(a)-\kappa_A(a)\|=0,
\]
since $k^{[0]}+k^{[1]}=1$. Choose one coordinate in the intersection
of the finitely many $\omega$-large approximation sets for
$a\in\mathcal F$. This proves the assertion. The first map into
$F_n\oplus F_n$ is $a\mapsto(\alpha_n(a),\alpha_n(a))$;
it is contractive because the direct-sum norm is the maximum norm.
\end{proof}

\begin{remark}[Keeping two colours under stabilization]
To pass from the corner inclusion to the algebra itself, we first
amplify by finite matrices.  Matrix amplification preserves each
order-zero return map, and the stabilization accommodates the second
diagonal block of the amplified target.  The resulting approximations
still have two colours.
\end{remark}

\begin{lemma}[Passing from a corner to the stabilization]
\label{end:stabilize}
Let $A$ be unital. Suppose that the first-corner inclusion
$A\to M_2(A)$ has finite-dimensional approximations with two colours,
with a c.p.c. first map and c.p.c. order-zero return maps.
Then $\nucdim(A\otimes\mathcal K)\le1$.
\end{lemma}

\begin{proof}
Fix $m\ge1$. For a finite set of elements of $A\otimes M_m$,
choose a two-colour approximation of $A\to M_2(A)$ accurate on
all their matrix entries. Amplify every map by $\id_{M_m}$.
The intermediate algebra is still the direct sum of two
finite-dimensional algebras, namely $E_0\otimes M_m$ and
$E_1\otimes M_m$. The first map remains c.p.c. Each amplified
return map remains c.p.c. order zero: if $\varphi=h\varphi_{\mathrm{hom}}$,
then
\[
 (\varphi\otimes\id_{M_m})(x)
 =(h\otimes1)(\varphi_{\mathrm{hom}}\otimes\id_{M_m})(x),
\]
with a commuting positive contractive weight.

Embed $M_2(A)\otimes M_m$ into $A\otimes\mathcal K$ so that its
first diagonal $m\times m$ block is the prescribed first $m$-corner.
The amplified approximation then approaches the inclusion of
$A\otimes M_m$ into $A\otimes\mathcal K$. Entrywise accuracy
suffices: an $m\times m$ matrix with entries of norm at most $\delta$
has norm at most $m\delta$.

Let $p_m=1_A\otimes\sum_{r=1}^m e_{rr}$. Precompose the amplified
first map with the c.p.c. compression $x\mapsto p_mxp_m$ on
$A\otimes\mathcal K$. For an arbitrary prescribed finite set in
the stabilization, first choose $m$ so that these compressions
approximate every element, then approximate the compressed finite
set as above. The resulting composite has the same value on
$x$ and $p_mxp_m$, so the total error is bounded by the compression
error plus the finite-corner approximation error. This gives the
required two-colour approximation on the stabilization.
\end{proof}

\begin{proof}[Proof of Theorem~\ref{intro:main}]
First let $A_0$ be nonzero and unital, and set $A=A_0\otimes\Z$.
Strong self-absorption gives isomorphisms
\[
 \sigma_l:\Z\otimes\Z\longrightarrow\Z,
 \qquad \sigma_l(z\otimes1)\longrightarrow z.
\]
Identify the three extra copies of $\Z$ with one copy by a fixed
unital isomorphism. The maps $\id_{A_0}\otimes\sigma_l$ then give
isomorphisms
\[
 \Gamma_l:A\otimes\Z^{\otimes3}\longrightarrow A,
 \qquad \Gamma_l(a\otimes1)\longrightarrow a\quad(a\in A).
\]
The convergence holds first on elementary tensors and then on all
of $A$ by contractivity.

Given a finite subset of $A$, choose $l$ so that $\Gamma_l$ has the
required accuracy on its first-factor images. Apply
Proposition~\ref{end:corner} with a sufficiently small error, and
postcompose its return maps with $\id_{M_2}\otimes\Gamma_l$.
This gives two-colour approximations to $A\to M_2(A)$.
Lemma~\ref{end:stabilize} shows that
$\nucdim(A\otimes\mathcal K)\le1$.

Nuclear dimension decreases on hereditary subalgebras
\cite[Proposition~2.5]{WZNuclear}. Applying this to the first corner
$A\otimes e_{11}$ gives $\nucdim(A_0\otimes\Z)\le1$.
For nonunital separable nuclear $A_0$, its unitization
$\widetilde A_0$ is separable and nuclear. The unital case applies
to $\widetilde A_0\otimes\Z$, and $A_0\otimes\Z$ is a closed
two-sided ideal, hence a hereditary subalgebra, of that algebra.
The same permanence result gives the bound. The zero algebra is
immediate.

Finally, if $A_0\cong A_0\otimes\Z$, transport these
finite-dimensional approximations across the isomorphism. The first
map is c.p.c., the two return maps are individually c.p.c. order zero,
and their sum approximates the identity, exactly as in
Definition~\ref{prelim:nuclear-dimension}.
\end{proof}

\section{A real tensor with uniform sign separation}
\label{sec:gaussian}
The dimension bound is now proved. We turn to the independent
splitting theorem that will supply absorption from finite nuclear
dimension. Its first ingredient is purely real and finite-dimensional:
one tensor must take both signs on every sufficiently separated product
of two-point sets. The following lemma obtains this simultaneously over
all such configurations. We equip real tensor powers with the Euclidean
tensor product inner product.

\begin{lemma}\label{lem:gaussian}
Fix \(k\geq2\) and \(0<r<1\). There are \(n\geq1\) and \(G\in(\R^k)^{\otimes n}\) such that, for every list of pairs of unit vectors \((u_j^0,u_j^1)_{j=1}^n\) satisfying
\[
 |\langle u_j^0,u_j^1\rangle|\leq r\qquad(1\leq j\leq n),
\]
the numbers
\[
 \big\langle G,u_1^{\alpha_1}\otimes\cdots\otimes u_n^{\alpha_n}\big\rangle,
 \qquad \alpha\in\{0,1\}^n,
\]
include a value greater than \(1/2\) and a value less than \(-1/2\).
\end{lemma}
\begin{proof}
Choose
\[
 0<a<\min\left\{\log2-\log(4/3)-\tfrac14\log3,\ -\tfrac14\log r\right\}.
\]
The first expression is positive. A binary Hamming ball of radius \(\lfloor n/4\rfloor\) has cardinality at most
\[
 \sum_{b=0}^{\lfloor n/4\rfloor}\binom nb
 \leq 3^{n/4}\sum_{b=0}^n\binom nb3^{-b}
 =3^{n/4}(4/3)^n.
\]
Gilbert's greedy deletion argument \cite[Theorem~1]{Gilbert}
therefore produces, for all sufficiently large \(n\), a set \(S_n\subset\{0,1\}^n\) of cardinality \(N=\lfloor e^{an}\rfloor\) whose distinct words have Hamming distance at least \(n/4\).

Let the coordinates of \(G\) be independent standard real Gaussian variables. Fix an admissible list and set
\[
 X_\alpha=\langle G,u_1^{\alpha_1}\otimes\cdots\otimes u_n^{\alpha_n}\rangle
 \quad(\alpha\in S_n).
\]
These variables are jointly centered Gaussian. Their covariance matrix \(\Sigma\) has diagonal entries one and off-diagonal entries of absolute value at most \(r^{n/4}\). Consequently its eigenvalues lie in \([1-\lambda_n,1+\lambda_n]\), where
\[
 \lambda_n=N r^{n/4}\longrightarrow0.
\]
For large \(n\), the density of this Gaussian vector is bounded pointwise by
\[
 \left(\frac{1+\lambda_n}{1-\lambda_n}\right)^{N/2}
\]
times the density of \(N\) independent centered Gaussians of variance \(1+\lambda_n\). Indeed, \(\det\Sigma\geq(1-\lambda_n)^N\) and \(\Sigma^{-1}\geq(1+\lambda_n)^{-1}I\).

Writing \(\Phi\) for the standard normal distribution function gives
\begin{equation}\label{eq:prob}
 \Pr(X_\alpha\leq1\ \forall\alpha\in S_n)
 \leq\left[\sqrt{\frac{1+\lambda_n}{1-\lambda_n}}\,
 \Phi\big((1+\lambda_n)^{-1/2}\big)\right]^N
 \leq q^N
\end{equation}
for one \(q<1\) and every sufficiently large \(n\). The last inequality is uniform over admissible lists because the bracket tends to \(\Phi(1)<1\). By symmetry the same bound holds for the event \(X_\alpha\geq-1\) for all \(\alpha\).

Set \(\eta_n=(4n k^{n/2})^{-1}\). The compact space of admissible lists has an \(\eta_n\)-net, in the maximum Euclidean distance between corresponding vectors, of size at most
\[
 (C_k/\eta_n)^{2kn}.
\]
For example, mesh the coordinate cubes and select an admissible representative from each occupied box. The logarithm of this bound is \(O(n^2)\). A union bound using \eqref{eq:prob} shows that the probability that some net list lacks a value above \(1\) or below \(-1\) is at most
\[
 2(C_k/\eta_n)^{2kn}q^N\longrightarrow0.
\]
Also \(\mathbb E\|G\|^2=k^n\), so
\(\Pr(\|G\|\leq2k^{n/2})\geq3/4\).
For large \(n\), both the successful-net event and this norm bound therefore hold for some \(G\).

If an admissible list is approximated by a net list, telescoping the elementary tensors of unit vectors bounds the change of each evaluation by
\[
 \|G\|\,n\eta_n\leq\tfrac12.
\]
The two evaluations exceeding \(1\) and below \(-1\) on the net list retain values above \(1/2\) and below \(-1/2\), respectively. This proves the assertion for every admissible list.
\end{proof}


\section{Testing signs in maximal tensor products}
\label{sec:tensor}
The goal is to prove Theorem~\ref{thm:tensor-intro}: the two signs of
one selfadjoint element will survive in every irreducible representation.
Such representations of a maximal tensor product need not be spatial.
The following factor argument and the pure-state excision method of
Akemann--Anderson--Pedersen \cite[Proposition~2.2]{AAP} let us transfer
the real evaluations of Lemma~\ref{lem:gaussian} without a spatial
representation hypothesis. We include the required excision proof.

\begin{lemma}[Products in commuting factors]\label{lem:factors}
Let \(M_1,\ldots,M_n\subseteq\mathcal B(H)\) be mutually commuting unital von Neumann factors. If \(e_j\in M_j\) is positive and \(\|e_j\|=1\), then
\[
 \left\|\prod_{j=1}^ne_j\right\|=1.
\]
\end{lemma}
\begin{proof}
First, a product of nonzero projections \(p_j\in M_j\) is nonzero. Induct on \(n\). By induction \(q=p_2\cdots p_n\) is a nonzero projection commuting with \(M_1\). The restriction \(M_1\to\mathcal B(qH)\) is a nonzero normal representation. Its kernel is a weak-* closed ideal in a factor and is therefore zero. Thus \(p_1q\ne0\).

For \(0<\varepsilon<1\), the spectral projections
\(p_j=\mathbf1_{(1-\varepsilon,1]}(e_j)\) are nonzero. Their product is nonzero, and on its range the positive product \(\prod_j e_j\) is at least \((1-\varepsilon)^n\). Letting \(\varepsilon\downarrow0\) proves the lower bound one; the reverse bound follows from contractivity.
\end{proof}

\begin{lemma}[Pure-state excision]\label{lem:excision}
Let \(E\) be a unital \(C^*\)-algebra and \(\rho\) a pure state. Given a finite set \(\mathcal F\subset E\) and \(\varepsilon>0\), there is a positive contraction \(f\in E\) with \(\|f\|=1\) such that
\[
 \|fxf-\rho(x)f^2\|<\varepsilon\qquad(x\in\mathcal F).
\]
\end{lemma}
\begin{proof}
Let \((\gamma,H,\xi)\) be the irreducible GNS representation of \(\rho\), with \(\|\xi\|=1\), and put
\[
 L=\{z\in E:\gamma(z)\xi=0\}.
\]
We first show that \(\ker\rho=L+L^*\). The inclusion from right to left is immediate. If \(\rho(x)=0\), write \(\eta=\gamma(x)\xi\perp\xi\). If \(\eta=0\), then \(x\in L\). Otherwise, on \(\operatorname{span}\{\xi,\eta\}\) let \(T\) be the rank-one operator with \(T\xi=\eta\) and \(T\eta=0\); thus \(T^*\xi=0\). Apply the selfadjoint form of Kadison transitivity \cite[Theorem~1]{Kadison} separately to the real and imaginary parts of \(T\). This gives \(y\in E\) such that
\[
 \gamma(y)\xi=\eta,\qquad\gamma(y^*)\xi=0.
\]
Then \(x-y\in L\) and \(y\in L^*\), as desired.

The algebra \(L\cap L^*\) is hereditary. Let \((g_\lambda)\) be a positive contractive approximate identity for it. If \(z\in L\), then \(z^*z\in L\cap L^*\), whence
\[
 \|z(1-g_\lambda)\|^2
 =\|(1-g_\lambda)z^*z(1-g_\lambda)\|\longrightarrow0.
\]
The adjoint assertion holds for elements of \(L^*\). Therefore, with \(f_\lambda=1-g_\lambda\), the equality \(\ker\rho=L+L^*\) gives
\[
 \|f_\lambda(x-\rho(x)1)f_\lambda\|\longrightarrow0
 \qquad(x\in E).
\]
Each \(f_\lambda\) is a positive contraction of norm one because \(\gamma(f_\lambda)\xi=\xi\). One sufficiently large \(\lambda\) works for \(\mathcal F\).
\end{proof}

\begin{proof}[Proof of Theorem~\ref{thm:tensor-intro}]
The closed ideal generated by commutators is all of \(D\); otherwise a nonzero unital commutative quotient would supply a character. Decomposing elements into selfadjoint parts and rescaling, choose selfadjoint contractions \(s_i,t_i\) and elements \(v_i,w_i\), for \(1\leq i\leq m\), such that
\begin{equation}\label{eq:commutator}
 \left\|1-\sum_{i=1}^m v_i[s_i,t_i]w_i\right\|<\tfrac12.
\end{equation}
Put \(s_0=1\), \(k=m+1\), and
\(M=\sum_i\|v_i\|\|w_i\|>0\).
Under any nonzero unital representation \(\sigma\), some represented commutator has norm greater than \(\delta=(2M)^{-1}\). Since
\[
 \|[S,T]\|\leq2\|S-c1\|\|T\|
 \leq\operatorname{diam}(\spc S)
\]
when \(S=S^*\), \(\|T\|\leq1\), and \(c\) is the midpoint of \(\spc S\), the spectrum of some \(\sigma(s_i)\) has diameter at least \(\delta\).

There are pure states \(\rho^0,\rho^1\) of \(\sigma(D)\) taking that element to its two spectral endpoints. Indeed, a state attaining an endpoint extends from the unital abelian algebra \(C^*(1,\sigma(s_i))\), and the states attaining that endpoint form a nonempty weak-* compact face of the state space. An extreme point of this face is pure. The associated real expectation vectors
\[
 v^\alpha=\big(\rho^\alpha(\sigma(s_b))\big)_{b=0}^m
 \in\{1\}\times[-1,1]^m\qquad(\alpha=0,1)
\]
differ by at least \(\delta\) in one coordinate. The set of pairs in \(\big(\{1\}\times[-1,1]^m\big)^2\) with \(\|v^0-v^1\|_\infty\geq\delta\) is compact. Consequently there is \(0<r<1\), independent of \(\sigma\), such that
\begin{equation}\label{eq:angle}
 \left|\left\langle\frac{v^0}{\|v^0\|},\frac{v^1}{\|v^1\|}\right\rangle\right|\leq r.
\end{equation}
To justify the strict bound, two vectors in this compact set cannot be negatively collinear because both first coordinates are one; positive collinearity forces them to coincide, contrary to their separation.

Choose \(G,n\) from Lemma~\ref{lem:gaussian} for \(k,r\) and form the selfadjoint element
\begin{equation}\label{eq:h}
 h=\sum_{b_1,\ldots,b_n=0}^m G_{b_1\ldots b_n}
 s_{b_1}\otimes\cdots\otimes s_{b_n}
 \in D^{\otimes_{\max}n}.
\end{equation}
Let \(\pi\) be any nonzero irreducible representation of this tensor product, with factor restrictions \(\pi_j:D\to\mathcal B(H)\). For each \(j\), choose the two pure states of \(\pi_j(D)\) described above. By \eqref{eq:angle} and Lemma~\ref{lem:gaussian}, choosing one state \(\rho_j\) from each pair makes
\begin{equation}\label{eq:p}
 p=\sum_{b_1,\ldots,b_n=0}^mG_{b_1\ldots b_n}
 \prod_{j=1}^n\rho_j(\pi_j(s_{b_j}))
\end{equation}
strictly positive for one choice and strictly negative for another. Removing the normalizations multiplies each tensor evaluation by a positive number and so preserves its sign.

The von Neumann algebras \(M_j=\pi_j(D)''\) are mutually commuting factors: a central element of any \(M_j\) commutes with every factor image and hence with the irreducible joint representation, so is scalar. For one of the choices in \eqref{eq:p}, use Lemma~\ref{lem:excision} to choose positive norm-one elements \(e_j\in\pi_j(D)\) with
\[
 \|e_j\pi_j(s_b)e_j-\rho_j(\pi_j(s_b))e_j^2\|<\varepsilon
 \quad(0\leq b\leq m).
\]
The product \(e=\prod_j e_j\) is positive and has norm one by Lemma~\ref{lem:factors}. Commuting distinct factors and telescoping their products in \eqref{eq:h} yields
\begin{equation}\label{eq:compression}
 \|e\pi(h)e-pe^2\|
 \leq n\varepsilon\sum_{b_1,\ldots,b_n}|G_{b_1\ldots b_n}|.
\end{equation}
All factors in this estimate are contractions.

If \(p>0\) and \(\pi(h)\leq0\), then \(pe^2-e\pi(h)e\geq pe^2\), so the left side of \eqref{eq:compression} is at least \(p\). Choosing \(\varepsilon\) sufficiently small is a contradiction. The negative choice of \(p\), with its own excisers, similarly rules out \(\pi(h)\geq0\). Thus
\[
 \pi(h_+)\ne0,\qquad\pi((-h)_+)\ne0
\]
for every nonzero irreducible \(\pi\). Each positive part is full: a proper closed ideal containing it would yield a nonzero quotient and hence an irreducible representation killing it. The two positive parts are orthogonal by functional calculus.
\end{proof}

The canonical unital map from the finite tensor power into the infinite maximal tensor power preserves fullness and orthogonality. Thus Theorem~\ref{thm:tensor-intro} also gives the affirmative answer to \cite[Question 6.2]{KR}.

\section{Central sequence algebras}\label{sec:central}
We now transfer the full signs of Theorem~\ref{thm:tensor-intro} to the
central sequence algebra and apply the absorption criterion of Robert
and Tikuisis. This completes the implication from finite nuclear dimension
to absorption.

Retain the free ultrafilter \(\omega\) on \(\N\). For a \(C^*\)-algebra \(A\), let
\[
 A_\omega=\ell^\infty(A)/\{(a_l):\lim_{l\to\omega}\|a_l\|=0\},\qquad
 F(A)=(A_\omega\cap A')/\ann(A),
\]
where \(A\) is embedded diagonally and
\[
 \ann(A)=\{b\in A_\omega:ba=ab=0\text{ for every }a\in A\}.
\]
For nonzero separable \(A\), this quotient is nonzero and unital. If \((e_l)\) is a sequential positive contractive approximate identity, its class is central and represents the unit. Indeed, for a bounded central sequence \((b_l)\) and \(a\in A\), commuting \(b_l\) past \(a\) shows \((e_lb_l-b_l)a\to_\omega0\), and likewise on the other side. Also \(e_la\to a\), so the unit is nonzero.

We use the following two results.
\begin{theorem}[Winter--Zacharias and Robert--Tikuisis]\label{thm:imports}
Let \(A\ne0\) be separable with \(\nucdim A\leq d<\infty\).
\begin{enumerate}
\item If no hereditary subalgebra of \(A\) has a finite-dimensional
irreducible representation, then there are finitely many finite-dimensional
algebras $E_j$, completely positive order-zero maps
$\varphi_j:E_j\to F(A)$, and elements $c_j\in F(A)$ such that every
matrix summand of every $E_j$ has size at least two and
\[
 1_{F(A)}=\sum_jc_j^*\varphi_j(1_{E_j})c_j.
\]
This is the matrix-threshold-two consequence of the covering-number
bound $\operatorname{cov}(F(A))\le(d+1)^2$
\cite[Definition~4.1 and Proposition~4.3]{WZNuclear}.
\item If \(F(A)\) contains two full orthogonal positive elements, then \(A\cong A\otimes_{\min}\Z\) \cite[Theorem 1.2]{RT}.
\end{enumerate}
Both assertions apply to nonunital \(A\), with the annihilator quotient just defined.
\end{theorem}

\begin{lemma}\label{lem:nocharacters}
If \(A\ne0\) is separable, nowhere scattered, and has finite nuclear dimension, then \(F(A)\) has no characters.
\end{lemma}
\begin{proof}
A hereditary subalgebra \(H\) is strongly Morita equivalent to the ideal \(I\) it generates. If \(H\) had a matrix quotient, ideal correspondence would give an ideal \(J\subseteq I\) for which \(I/J\) is strongly Morita equivalent to a matrix algebra, hence to \(\C\). By imprimitivity, \(I/J\) would be elementary. Ideals of \(I\) are ideals of \(A\), contradicting nowhere scatteredness. This is also the equivalence in \cite[Theorem 3.1]{TV}.

Apply Theorem~\ref{thm:imports}(1) to $A$, put $B=F(A)$, and retain its
maps $\varphi_j:E_j\to B$, elements $c_j$, and displayed identity.

If \(\chi:B\to\C\) were a character, its compositions with these maps would be order-zero positive functionals. Such a functional on \(M_r\), \(r\geq2\), is zero. To see this, a nonzero positive functional has a positive density matrix with an eigenvector \(\xi\) of positive eigenvalue. For a unit eigenvector \(\eta\perp\xi\), the orthogonal rank-one projections onto \((\xi+\eta)/\sqrt2\) and \((\xi-\eta)/\sqrt2\) both have strictly positive values, contrary to order zero. Thus \(\chi\circ\varphi_j=0\) for every \(j\), contradicting the displayed identity and \(\chi(1_B)=1\).
\end{proof}

\begin{lemma}[Commuting reindexed copies]\label{lem:reindex}
Let \(A\ne0\) be separable, and let \(D\subseteq F(A)\) be a separable unital subalgebra with the same unit. For every finite \(n\), there is a unital \(*\)-homomorphism
\[
 D^{\otimes_{\max}n}\longrightarrow F(A).
\]
\end{lemma}
\begin{proof}
It suffices to prove that for every separable subalgebra \(C\subseteq F(A)\), there is a unital \(*\)-homomorphism from \(D\) into \(F(A)\) with range commuting with \(C\). This is Kirchberg's reindexing mechanism, in the form recorded in
\cite[Proposition~2.2 and Corollary~2.3]{KR} for unital $A$. We include
the annihilator details needed for nonunital $A$.

Choose a countable dense unital \(*\)-algebra \(D_0\subset D\) over \(\Q+i\Q\). For each \(x\in D_0\), choose a bounded lift \((a_l^x)_l\) in \(A\) that represents a central element mapping to \(x\), with
\begin{equation}\label{eq:liftbound}
 \sup_l\|a_l^x\|\leq2\|x\|.
\end{equation}
Take the zero lift for zero and the approximate identity for \(1\). The quotient norm permits these choices: first lift within a bound strictly below \(2\|x\|\), then truncate coordinate norms without changing the ultrapower class. The additive, scalar, multiplicative, and adjoint relation defects between these lifts annihilate \(A\) on both sides along \(\omega\).

Choose a countable dense family in \(C\) with bounded lifts \((b_s^u)_s\) in \(A\). Also fix a countable dense subset of \(A\). List the following countably many requirements on an old index \(l\): centrality of each \(a_l^x\) against that subset; left and right products of every relation defect with that subset tending to zero; and
\[
 a_l^1a\longrightarrow a,\qquad aa_l^1\longrightarrow a
\]
for its elements \(a\). All these requirements hold along \(\omega\).

For each output index \(s\), choose \(l(s)\) so that the first \(s\) requirements hold with error less than \(1/s\), and also
\[
 \|[a_{l(s)}^x,b_s^u]\|<1/s
\]
for the first \(s\) pairs \((x,u)\) in an enumeration. For this fixed \(s\), the finitely many \(b_s^u\) are fixed elements of \(A\). Original centrality therefore makes all these commutator conditions \(\omega\)-large as conditions on \(l\). A finite intersection with the other \(\omega\)-large conditions is nonempty, so the choice is possible.

The new sequences \((a_{l(s)}^x)_s\) are central, and all their relation defects annihilate \(A\) on both sides: the tests hold in the ordinary limit, and the uniform bounds extend them from the dense subset to all of \(A\). Thus they define a unital rational-complex \(*\)-homomorphism on \(D_0\) into \(F(A)\). The unit assertion follows from the approximate-identity tests. Bound \eqref{eq:liftbound} permits continuous extension to a complex unital \(*\)-homomorphism on \(D\). Its range commutes with the dense family in \(C\), hence with \(C\). Injectivity is not needed.

Starting with one copy and repeatedly taking \(C\) to be generated by the previous images gives \(n\) commuting unital images of \(D\). The universal property of the maximal tensor product gives the stated homomorphism.
\end{proof}

\begin{proof}[Proof of Theorem~\ref{thm:main}]
The assertion is immediate for \(A=0\), so assume \(A\ne0\).
Assume first that $\nucdim(A)<\infty$; we prove (1)$\Rightarrow$(2).
By Lemma~\ref{lem:nocharacters}, \(B=F(A)\) has no characters. Choose in \(B\) a finite commutator witness as in \eqref{eq:commutator}, and let \(D\) be the unital \(C^*\)-subalgebra generated by all its entries. Then \(D\) is separable and has no characters, since applying a character of \(D\) to that inequality would give \(1<1/2\).

By Theorem~\ref{thm:tensor-intro}, some \(D^{\otimes_{\max}n}\) contains full orthogonal positive elements \(a,b\). Lemma~\ref{lem:reindex} supplies a unital homomorphism into \(B\). Their images remain orthogonal and full: applying a unital homomorphism to finite ideal combinations approximating the unit proves that fullness passes to the target. Theorem~\ref{thm:imports}(2) now gives \(A\cong A\otimes_{\min}\Z\).
Conversely, for separable nuclear $A$, absorption and
Theorem~\ref{intro:main} imply $\nucdim(A)\le1$. This bound implies
finite nuclear dimension, completing all three equivalences in
Theorem~\ref{thm:main}.
\end{proof}

\begin{remark}
The same argument proves that \(F(A)\) has two full orthogonal positive elements whenever \(A\ne0\) is separable and \(F(A)\) has no characters. Finite nuclear dimension enters the proof of Theorem~\ref{thm:main} only through the two imported results in Theorem~\ref{thm:imports}.
\end{remark}

\subsection{Cuntz-semigroup consequence}
\label{subsec:cu-consequence}

We use the completed stabilized Cuntz semigroup $\Cu(X)$ formed from
$(X\otimes\K)_+$, with no distinguished scale. Isomorphisms in
$\mathbf{Cu}$ preserve zero, addition, order, increasing sequential
suprema, and the way-below relation in both directions. The isomorphism
in Corollary~\ref{cor:regularity-cu}\textup{(ii)} need not agree with
$\Cu(\iota_A)$.

\begin{proof}[Proof of Corollary~\ref{cor:regularity-cu}]
Theorem~\ref{thm:main} gives the equivalence of its three conditions.
The companion \cite[Corollaries~1.3 and~1.4]{CuntzCompanion} gives
that \textup{(i)} implies absorption and that \textup{(ii)} is equivalent
to absorption, even without the exclusion on elementary ideal subquotients.

\Needspace{7\baselineskip}
To prove that absorption implies \textup{(i)}, the case $A=0$ is immediate.
Put $D=A\otimes\Z$ and $j_X(x)=x\otimes1_\Z$, so $j_A=\iota_A$.
Strong self-absorption gives an isomorphism
$\varphi:\Z\to\Z\otimes\Z$ approximately unitarily equivalent to
$z\mapsto z\otimes1_\Z$. The companion's nonunital calculation
\cite[proof of Corollary~1.4]{CuntzCompanion} gives
$\Cu(j_D)=\Cu(\id_A\otimes\varphi)$, an isomorphism: the implementing
unitaries $u_n\in\Z\otimes\Z$ act through
$1_{M(A)}\otimes u_n\in M(A\otimes\Z\otimes\Z)$, and positive cuts
and approximate units after stabilization give equality of the completed
Cuntz maps.

For an absorption isomorphism $\theta:A\to D$, both sides of the first
identity below send $a$ to $\theta(a)\otimes1_\Z$. Functoriality then gives
\[
\begin{aligned}
 (\theta\otimes\id_\Z)\circ j_A&=j_D\circ\theta,\\
 \Cu(j_A)&=\Cu(\theta\otimes\id_\Z)^{-1}
       \circ\Cu(j_D)\circ\Cu(\theta).
\end{aligned}
\]
All three factors on the right are isomorphisms, proving \textup{(i)}.
Thus \textup{(ii)} reaches the canonical condition \textup{(i)} through
absorption, not by identifying its abstract isomorphism with $\Cu(\iota_A)$.
\end{proof}

\begin{thebibliography}{BBSTWW}
\raggedright

\bibitem[AAP]{AAP}
Charles A. Akemann, Joel Anderson, and Gert K. Pedersen,
\emph{Excising states of $C^*$-algebras},
Canad. J. Math. \textbf{38} (1986), no.~5, 1239--1260.
\href{https://doi.org/10.4153/CJM-1986-063-7}{doi:10.4153/CJM-1986-063-7}.

\bibitem[APTV]{APTV}
Ramon Antoine, Francesc Perera, Hannes Thiel, and Eduard Vilalta,
\emph{Pure $C^*$-algebras},
preprint, \href{https://arxiv.org/abs/2406.11052v3}{arXiv:2406.11052v3},
2024.

\bibitem[BBSTWW]{BBSTWW}
Joan Bosa, Nathanial P. Brown, Yasuhiko Sato, Aaron Tikuisis,
Stuart White, and Wilhelm Winter,
\emph{Covering dimension of $C^*$-algebras and 2-coloured classification},
Mem. Amer. Math. Soc. \textbf{257} (2019), no.~1233, vii+97 pp.
\href{https://doi.org/10.1090/memo/1233}{doi:10.1090/memo/1233}.
\href{https://arxiv.org/abs/1506.03974v3}{arXiv:1506.03974v3}.

\bibitem[BGSW]{BGSW}
Joan Bosa, James Gabe, Aidan Sims, and Stuart White,
\emph{The nuclear dimension of $\mathcal O_\infty$-stable $C^*$-algebras},
Adv. Math. \textbf{401} (2022), 108250.
\href{https://doi.org/10.1016/j.aim.2022.108250}{doi:10.1016/j.aim.2022.108250}.
\href{https://arxiv.org/abs/1906.02066}{arXiv:1906.02066}.

\bibitem[BCW]{BCW}
Nathanial P. Brown, Jos\'e R. Carri\'on, and Stuart White,
\emph{Decomposable approximations revisited},
in \emph{Operator Algebras and Applications: The Abel Symposium 2015}
(Toke M. Carlsen, Nadia S. Larsen, Sergey Neshveyev, and Christian Skau, eds.),
Abel Symposia, vol.~12, Springer, Cham, 2016, pp.~45--65.
\href{https://doi.org/10.1007/978-3-319-39286-8_2}{doi:10.1007/978-3-319-39286-8\_2}.
\href{https://arxiv.org/abs/1602.00021v4}{arXiv:1602.00021v4}.

\bibitem[CE]{CE}
Jorge Castillejos and Samuel Evington,
\emph{Nuclear dimension of simple stably projectionless $C^*$-algebras},
Anal. PDE \textbf{13} (2020), no.~7, 2205--2240.
\href{https://doi.org/10.2140/apde.2020.13.2205}{doi:10.2140/apde.2020.13.2205}.
\href{https://arxiv.org/abs/1901.11441v4}{arXiv:1901.11441v4}.

\bibitem[CETWW]{CETWW}
Jorge Castillejos, Samuel Evington, Aaron Tikuisis, Stuart White,
and Wilhelm Winter,
\emph{Nuclear dimension of simple $C^*$-algebras},
Invent. Math. \textbf{224} (2021), no.~1, 245--290.
\href{https://doi.org/10.1007/s00222-020-01013-1}{doi:10.1007/s00222-020-01013-1}.
\href{https://arxiv.org/abs/1901.05853v3}{arXiv:1901.05853v3}.

\bibitem[ChE]{ChoiEffros}
Man-Duen Choi and Edward G. Effros,
\emph{The completely positive lifting problem for $C^*$-algebras},
Ann. of Math. (2) \textbf{104} (1976), no.~3, 585--609.
\href{https://doi.org/10.2307/1970968}{doi:10.2307/1970968}.

\bibitem[ENST]{ENST}
George A. Elliott, Zhuang Niu, Luis Santiago, and Aaron Tikuisis,
\emph{Decomposition rank of approximately subhomogeneous $C^*$-algebras},
Forum Math. \textbf{32} (2020), no.~4, 827--889.
\href{https://doi.org/10.1515/forum-2020-0018}{doi:10.1515/forum-2020-0018}.
\href{https://arxiv.org/abs/1505.06100}{arXiv:1505.06100}.

\bibitem[ERS]{ERS}
George A. Elliott, Leonel Robert, and Luis Santiago,
\emph{The cone of lower semicontinuous traces on a $C^*$-algebra},
Amer. J. Math. \textbf{133} (2011), no.~4, 969--1005.
\href{https://doi.org/10.1353/ajm.2011.0027}{doi:10.1353/ajm.2011.0027}.
\href{https://arxiv.org/abs/0805.3122v2}{arXiv:0805.3122v2}.

\bibitem[Gabe]{Gabe}
James Gabe,
\emph{A new proof of Kirchberg's $\mathcal O_2$-stable classification},
J. Reine Angew. Math. \textbf{761} (2020), 247--289.
\href{https://doi.org/10.1515/crelle-2018-0010}{doi:10.1515/crelle-2018-0010}.
\href{https://arxiv.org/abs/1706.03690v2}{arXiv:1706.03690v2}.

\bibitem[Gil]{Gilbert}
Edward N. Gilbert,
\emph{A comparison of signalling alphabets},
Bell System Tech. J. \textbf{31} (1952), no.~3, 504--522.
\href{https://doi.org/10.1002/j.1538-7305.1952.tb01393.x}{doi:10.1002/j.1538-7305.1952.tb01393.x}.

\bibitem[HKW]{HKW}
Ilan Hirshberg, Eberhard Kirchberg, and Stuart White,
\emph{Decomposable approximations of nuclear $C^*$-algebras},
Adv. Math. \textbf{230} (2012), no.~3, 1029--1039.
\href{https://doi.org/10.1016/j.aim.2012.03.028}{doi:10.1016/j.aim.2012.03.028}.
\href{https://arxiv.org/abs/1109.2379}{arXiv:1109.2379}.

\bibitem[JS]{JS}
Xinhui Jiang and Hongbing Su,
\emph{On a simple unital projectionless $C^*$-algebra},
Amer. J. Math. \textbf{121} (1999), no.~2, 359--413.
\href{https://doi.org/10.1353/ajm.1999.0012}{doi:10.1353/ajm.1999.0012}.

\bibitem[Kad]{Kadison} R.~V.~Kadison,
\emph{Irreducible operator algebras},
Proc. Natl. Acad. Sci. USA \textbf{43} (1957), 273--276.
\href{https://doi.org/10.1073/pnas.43.3.273}{doi:10.1073/pnas.43.3.273}.

\bibitem[KR05]{KR05}
Eberhard Kirchberg and Mikael R\o rdam,
\emph{Purely infinite $C^*$-algebras: ideal-preserving zero homotopies},
Geom. Funct. Anal. \textbf{15} (2005), no.~2, 377--415.
\href{https://arxiv.org/abs/math/0312286}{arXiv:math/0312286}.

\bibitem[KR]{KR} E.~Kirchberg and M.~R\o rdam,
\emph{When central sequence \(C^*\)-algebras have characters},
Int. J. Math. \textbf{26} (2015), no.~7, 1550049, 32 pp.
\href{https://doi.org/10.1142/S0129167X15500494}{doi:10.1142/S0129167X15500494}.

\bibitem[MS]{MS}
Hiroki Matui and Yasuhiko Sato,
\emph{Decomposition rank of UHF-absorbing $C^*$-algebras},
Duke Math. J. \textbf{163} (2014), no.~14, 2687--2708.
\href{https://doi.org/10.1215/00127094-2826908}{doi:10.1215/00127094-2826908}.
\href{https://arxiv.org/abs/1303.4371v2}{arXiv:1303.4371v2}.

\bibitem[OAI-Cu]{CuntzCompanion}
OpenAI,
\emph{Cuntz comparison and Jiang--Su absorption},
OpenAI Math Release preprint
\href{https://github.com/openai/math/blob/main/preprints/Cuntz-comparison-and-Jiang-Su-absorption-September-23-2026/paper.pdf}{OAI:Cuntz-comparison-and-Jiang-Su-absorption-September-23-2026},
2026.
Corollaries~1.3 and~1.4.

\bibitem[OAI]{Fixed324}
OpenAI,
\emph{The Kirchberg--R\o rdam character criterion},
OpenAI Math Release preprint
\href{https://github.com/openai/math/blob/main/preprints/The-Kirchberg-Rordam-character-criterion-September-25-2026/paper.pdf}{OAI:The-Kirchberg-Rordam-character-criterion-September-25-2026},
2026.
Theorems~1.1 and~1.2.

\bibitem[PTV]{PTV}
Francesc Perera, Hannes Thiel, and Eduard Vilalta,
\emph{Centrally pure $C^*$-algebras},
J. Funct. Anal. \textbf{291} (2026), 111578.
\href{https://doi.org/10.1016/j.jfa.2026.111578}{doi:10.1016/j.jfa.2026.111578}.
\href{https://arxiv.org/abs/2512.17261}{arXiv:2512.17261}.

\bibitem[RS]{RS}
Leonel Robert and Luis Santiago,
\emph{Classification of $C^*$-homomorphisms from $C_0(0,1]$ to a $C^*$-algebra},
J. Funct. Anal. \textbf{258} (2010), no.~3, 869--892.
\href{https://doi.org/10.1016/j.jfa.2009.06.025}{doi:10.1016/j.jfa.2009.06.025}.

\bibitem[RT]{RT} L.~Robert and A.~Tikuisis,
\emph{Nuclear dimension and \(\mathcal Z\)-stability of non-simple \(C^*\)-algebras},
Trans. Amer. Math. Soc. \textbf{369} (2017), no.~7, 4631--4670.
\href{https://doi.org/10.1090/tran/6842}{doi:10.1090/tran/6842}.
\href{https://arxiv.org/abs/1308.2941v2}{arXiv:1308.2941v2}.

\bibitem[R\o r]{Rordam}
Mikael R\o rdam,
\emph{The stable and the real rank of $\mathcal Z$-absorbing $C^*$-algebras},
Internat. J. Math. \textbf{15} (2004), no.~10, 1065--1084.
\href{https://doi.org/10.1142/S0129167X04002661}{doi:10.1142/S0129167X04002661}.
\href{https://arxiv.org/abs/math/0408020v2}{arXiv:math/0408020v2}.

\bibitem[SWW]{SWW}
Yasuhiko Sato, Stuart White, and Wilhelm Winter,
\emph{Nuclear dimension and $\mathcal Z$-stability},
Invent. Math. \textbf{202} (2015), no.~2, 893--921.
\href{https://doi.org/10.1007/s00222-015-0580-1}{doi:10.1007/s00222-015-0580-1}.
\href{https://arxiv.org/abs/1403.0747}{arXiv:1403.0747}.

\bibitem[STW]{STW99}
Christopher Schafhauser, Aaron Tikuisis, and Stuart White,
\emph{Nuclear $C^*$-algebras: 99 problems},
preprint, \href{https://arxiv.org/abs/2506.10902v2}{arXiv:2506.10902v2},
2025, revised 8 May 2026.

\bibitem[Sza]{SzaboOInfinity}
G\'abor Szab\'o,
\emph{On the nuclear dimension of strongly purely infinite $C^*$-algebras},
preprint, \href{https://arxiv.org/abs/1510.01917v2}{arXiv:1510.01917v2},
2015, revised 9 November 2016.

\bibitem[TV]{TV} H.~Thiel and E.~Vilalta,
\emph{Nowhere scattered \(C^*\)-algebras},
J. Noncommut. Geom. \textbf{18} (2024), 231--263.
\href{https://ems.press/content/serial-article-files/42631}{Publisher's text}.

\bibitem[Tik]{TikuisisFiniteIdeals}
Aaron Tikuisis,
\emph{Nuclear dimension, $\mathcal Z$-stability, and algebraic simplicity
for stably projectionless $C^*$-algebras},
preprint, \href{https://arxiv.org/abs/1210.2237v2}{arXiv:1210.2237v2},
2012, revised 29 May 2013.

\bibitem[TiW]{TiW}
Aaron Tikuisis and Wilhelm Winter,
\emph{Decomposition rank of $\mathcal Z$-stable $C^*$-algebras},
Anal. PDE \textbf{7} (2014), no.~3, 673--700.
\href{https://doi.org/10.2140/apde.2014.7.673}{doi:10.2140/apde.2014.7.673}.

\bibitem[TW]{TW}
Andrew S. Toms and Wilhelm Winter,
\emph{Strongly self-absorbing $C^*$-algebras},
Trans. Amer. Math. Soc. \textbf{359} (2007), no.~8, 3999--4029.
\href{https://doi.org/10.1090/S0002-9947-07-04173-6}{doi:10.1090/S0002-9947-07-04173-6}.

\bibitem[Win]{WinterAbsorption}
Wilhelm Winter,
\emph{Nuclear dimension and $\mathcal Z$-stability of pure $C^*$-algebras},
preprint, \href{https://arxiv.org/abs/1006.2731v2}{arXiv:1006.2731v2},
2010, revised 20 May 2011.

\bibitem[WZ1]{WZOrderZero}
Wilhelm Winter and Joachim Zacharias,
\emph{Completely positive maps of order zero},
M\"unster J. Math. \textbf{2} (2009), 311--324.
\href{https://www.uni-muenster.de/FB10/mjm/vol_2/mjm_vol_2_14.pdf}{Published version}.

\bibitem[WZ2]{WZNuclear}
Wilhelm Winter and Joachim Zacharias,
\emph{The nuclear dimension of $C^*$-algebras},
Adv. Math. \textbf{224} (2010), no.~2, 461--498.
\href{https://doi.org/10.1016/j.aim.2009.12.005}{doi:10.1016/j.aim.2009.12.005}.
\href{https://arxiv.org/abs/0903.4914v1}{arXiv:0903.4914v1}.

\end{thebibliography}

\end{document}
